\documentclass[11pt,letterpaper,reqno]{amsart}

\usepackage{tikz}
\usetikzlibrary{positioning, shapes.geometric, arrows.meta}
\usepackage{amssymb}
\usepackage{amsmath}
\usepackage{mathtools}
\usepackage{amsthm}
\usepackage{amsfonts}
\usepackage{bbm}
\usepackage{enumitem}
\usepackage{pgfplots}
\pgfplotsset{compat=1.18}
\usepackage{booktabs}
\usepackage{graphicx}
\usepackage[T1]{fontenc}
\usepackage{doi}
\usepackage{float}
\usepackage{microtype}
\addtolength{\hoffset}{-1.5cm}\addtolength{\textwidth}{3cm}
\addtolength{\voffset}{-1cm}\addtolength{\textheight}{2cm}
\usepackage{xcolor}
\usepackage{hyperref}
\usepackage{bookmark}
\hypersetup{
  pdfstartview={FitH},
  colorlinks=true,
  linkcolor=blue,
  citecolor=blue,
  urlcolor=blue
}
\emergencystretch=2em


\newtheorem{thm}{Theorem}[section]
\newtheorem{lem}[thm]{Lemma}
\newtheorem{prop}[thm]{Proposition}
\newtheorem{cor}[thm]{Corollary}
\newtheorem{claim}[thm]{Claim}
\newtheorem{prob}[thm]{Problem}
\newtheorem{ques}[thm]{Question}
\newtheorem{conj}[thm]{Conjecture}
\theoremstyle{definition}
\newtheorem{eg}[thm]{Example}
\newtheorem{rem}[thm]{Remark}
\newtheorem{defin}[thm]{Definition}
\numberwithin{equation}{section}

\newcommand{\YixinEmail}{yixin.he717@gmail.com} \newcommand{\QuanyuEmail}{tangquanyu827@gmail.com}
\newcommand{\sa}{\mathrm{sa}}
\newcommand{\RR}{\mathbb R}
\newcommand{\CC}{\mathbb C}
\newcommand{\NN}{\mathbb N}
\newcommand{\one}{1}
\newcommand{\ot}{\overline{\otimes}}
\newcommand{\cl}{\operatorname{cl}}
\DeclareMathOperator{\spanop}{span}
\DeclareMathOperator{\dist}{dist}
\DeclareMathOperator{\Ext}{Ext}
\DeclareMathOperator{\dens}{dens}
\DeclareMathOperator{\supp}{supp}
\DeclareMathOperator{\Hull}{Hull}

\begin{document}
\title[Unverified candidate proof for the nonseparable case] {An unverified AI-generated candidate proof for the nonseparable case of Kadison's trace-vector basis problem for type $\mathrm{II}_1$ factors}

\author{} 
\date{Public-review draft, July 24, 2026}


\begin{abstract}
In 1967, Kadison asked whether every type $\mathrm{II}_1$ factor has an
orthonormal basis, with respect to its trace, consisting of unitaries.
The case in which $L^2(M,\tau)$ is separable was recently settled
affirmatively.

This manuscript presents an AI-generated and independently unverified
candidate proof for the complementary nonseparable case in the tracial
standard representation.  More precisely, the candidate argument claims
that if $M$ is a type $\mathrm{II}_1$ factor with faithful normal tracial
state $\tau$ and
\[
   \dens L^2(M,\tau)>\aleph_0,
\]
then $L^2(M,\tau)$ admits a Hilbert orthonormal basis whose vectors are
represented by self-adjoint unitaries in $M$, and hence an orthonormal basis
consisting of trace vectors for the left action of $M$.

If correct, this argument, together with the previously established
separable case, would give an affirmative solution to Kadison's 1967 problem
in the tracial standard representation.  The candidate proof is based on a
relative norming lemma under small-density constraints, a finite-layer
certified-background formalism designed to keep Hilbert-space projections
represented by bounded operators in the ambient factor, and a transfinite
extension over the density character of $L^2(M,\tau)$.

The mathematical argument has not been independently verified.  This
manuscript is released solely to invite careful scrutiny, correction, and
possible further collaboration by specialists.
\end{abstract}





\maketitle



\begin{center} 
\begin{minipage}{0.92\textwidth} 
\small 
\textbf{Provenance and verification status.} This manuscript was generated by OpenAI's GPT-5.5 Pro in response to prompts and follow-up instructions supplied by Yixin He and Quanyu Tang. Yixin He and Quanyu Tang selected the problem, conducted the prompting process, compiled the generated material, and prepared it for public release. They are listed solely as project contacts and provenance contributors, not as authors who have verified the mathematical proof.


The argument has not been independently checked, and no claim is made that the stated theorem is correct. The theorem and proof environments below record the claims of the AI-generated manuscript rather than established results. Readers are strongly encouraged to report gaps, errors, or relevant prior work.



\medskip
\noindent
\textbf{Project contacts and affiliations:}

\medskip
\noindent
\textbf{Yixin He}\\
School of Mathematical Sciences, Fudan University,\\
Shanghai 200433, P.~R.~China.

\medskip
\noindent
\textbf{Quanyu Tang} \textit{(corresponding project contact)}\\
School of Mathematical Sciences, University of Science and Technology of China,\\
Hefei 230026, P.~R.~China.\\
\href{mailto:tangquanyu827@gmail.com}
{\nolinkurl{tangquanyu827@gmail.com}}
\end{minipage}
\end{center}





\bigskip





\section{Outline of the proof}\label{sec:outline}

This section records the structure of the proof and fixes the main pieces of
notation used below. The proof is organized around the real Hilbert
space
\[
   H_0(P)=L^2(P,\tau)_{\sa}\ominus \RR\one
\]
associated with a finite von Neumann algebra $P$.  By
Proposition~\ref{prop:trace-vectors-unitaries} and
Lemma~\ref{lem:trace-zero-reduction}, it is enough to construct an
orthonormal basis of $H_0(M)$ consisting of trace-zero symmetries in $M$.
Adding the vector $\one$ then gives a complex Hilbert-space orthonormal basis
of $L^2(M,\tau)$ represented by self-adjoint unitaries, and hence by trace
vectors.

The main new ingredient in the nonseparable case is the relative norming lemma,
Lemma~\ref{lem:large-relative-norming}.  Suppose $N\subset M$ has smaller
$L^2$-density than $M$, and suppose that $r=r^*\in M$ is orthogonal to
$L^2(N)_{\sa}$ and to a prescribed set $\Gamma$ of fewer than
$\dens L^2(M)$ self-adjoint constraints.  The lemma constructs a trace-zero
symmetry $s\in M$ which remains orthogonal to $N$ and to $\Gamma$, and which
satisfies the quantitative norming identity
\[
   \langle r,s\rangle_2=\frac{\|r\|_2^2}{\|r\|_\infty}.
\]
The proof uses a Krein--Milman extreme-point argument in a compact convex set
of positive contractions.  The only place where nonseparability enters is the
small-constraint perturbation lemma, which uses the strict density gap
$\dens L^2(N)<\dens L^2(M)$ to find a nonzero perturbation in a corner while
annihilating all prescribed constraints.

The second technical point is that, for an uncountable orthonormal family, the
Hilbert-space projection onto an arbitrary subfamily need not be represented by
an element of $M$.  Sections~\ref{sec:certified-backgrounds} and
\ref{sec:strong-closure} introduce finite-layer certified backgrounds to avoid
this problem.  A completed block
\[
   L^2(Q)_{\sa}=L^2(A)_{\sa}\oplus \overline{\spanop_{\RR}}T
\]
has the bounded projection formula
\[
   P_T(x)=E_Q(x)-E_A(E_Q(x)),\qquad x=x^*\in M.
\]
Finite-layer certificates are built by adjoining such blocks one layer at a
time.  The strong closure lemma, Proposition~\ref{prop:strong-closure}, says
that every small subfamily of a certified background can be enlarged to a still
small certified subbackground.  This is the bookkeeping device which guarantees
that all projections used later are bounded elements of $M$.

The countable extension lemma, Lemma~\ref{lem:countable-extension}, is the
basic construction step.  Starting from a small subalgebra $N$, a certified
background $\Omega$ orthogonal to $N$, and countable sets of new elements to
absorb, it constructs a separable completed block
\[
   L^2(R)_{\sa}=L^2(C)_{\sa}\oplus
   \overline{\spanop_{\RR}}(\Omega_0\cup U),
\]
where $C\subset N$, the set $\Omega_0\subset\Omega$ is closed and certified,
and $U$ is a countable orthonormal family of trace-zero symmetries orthogonal
to both $N$ and $\Omega$.  The proof repeatedly processes all eventual
self-adjoint $*$-polynomial words.  For a word $y$ one writes
\[
   b(y)=y-E_N(y)-P_\Omega(y).
\]
Whenever the residual part of $b(y)$ outside the already chosen span is
nonzero, Lemma~\ref{lem:large-relative-norming} supplies a new symmetry which
reduces the squared distance by
\[
   \frac{\|r\|_2^4}{\|r\|_\infty^2}.
\]
The bookkeeping is arranged so that each word is visited at times $n_k$ with
$\sum_k n_k^{-1}=\infty$; the elementary bound $\|r\|_\infty\le K_y\sqrt{n_k}$
then forces the residual distance to converge to zero.

Proposition~\ref{prop:LBE} upgrades the countable extension lemma to all
cardinals $\lambda<\dens L^2(M)$ by induction on $\lambda$.  The proof uses no
regularity assumption on the ambient density character; the needed cardinal
estimates are recorded in Remark~\ref{rem:no-regularity}.  At successor stages
one certifies the old background together with the old newly constructed layer
and applies the induction hypothesis again.  At limit stages one takes unions
of nested completed blocks and uses Lemma~\ref{lem:nested-limit}.

Finally, Proposition~\ref{prop:RE} converts the certified extension theorem
into a relative basis extension property: if $N\subset M$ has smaller density
and $H_0(N)$ already has a symmetry basis, then after adjoining any set of at
most that density many self-adjoint elements one obtains a larger subalgebra
$P$ whose $H_0(P)$ again has a symmetry basis extending the old one.  The main
nonseparable theorem is then a transfinite recursion over a dense family
$(a_\alpha)_{\alpha<\kappa}$ in $H_0(M)$, where
$\kappa=\dens L^2(M,\tau)$.  Starting from a separable diffuse abelian
subalgebra with its Walsh symmetry basis, the recursion builds nested
subalgebras $N_\alpha$ and nested symmetry bases $B_\alpha$, ensuring that
$a_\alpha\in N_{\alpha+1}$.  The union of the $B_\alpha$'s is therefore an
orthonormal basis of $H_0(M)$, and Lemma~\ref{lem:trace-zero-reduction}
completes the proof.







\section{Introduction and problem statement}\label{sec:introduction}

Let \(M\) be a finite von Neumann algebra equipped with a faithful normal
tracial state \(\tau\).  Throughout the paper, all traces are normalized
unless explicitly stated otherwise.  We denote by \(L^2(M,\tau)\) the Hilbert
space completion of \(M\) for the norm \(\|x\|_2=\tau(x^*x)^{1/2}\). For
\(x\in M\), its image in \(L^2(M,\tau)\) will be denoted by \(\widehat{x}\). We write
$
   M_{\mathrm{sa}}=\{x\in M:x=x^*\}
$
for the self-adjoint part of \(M\).  We also denote by
\(L^2(M,\tau)_{\mathrm{sa}}\) the real Hilbert space obtained as the
\(\|\cdot\|_2\)-closure of \(M_{\mathrm{sa}}\) in \(L^2(M,\tau)\). 
We use the convention that the inner product is linear in the first variable:
\[
   \langle \widehat{x},\widehat{y}\rangle_2=\tau(y^*x),
   \qquad x,y\in M.
\]
When no confusion is possible, we identify an element \(x\in M\) with its
vector \(\widehat x\in L^2(M,\tau)\).  Throughout the paper, \(M\) acts on
\(L^2(M,\tau)\) by left multiplication.  A vector
\(\xi\in L^2(M,\tau)\) is called a \emph{trace vector} for this action if
\[
   \langle x\xi,\xi\rangle_2=\tau(x),\qquad x\in M.
\]
By an \emph{orthonormal basis} we always mean a Hilbert-space orthonormal
basis, namely an orthonormal family whose closed linear span is the whole
Hilbert space.  In the nonseparable case such a basis may be uncountable.
In this paper, ``nonseparable'' refers to the Hilbert-space density character
\(\dens L^2(M,\tau)\), not to an arbitrary representation of \(M\).

We also denote by \(\mathcal U(M)\) the unitary group of \(M\):
\[
   \mathcal U(M)=\{u\in M:u^*u=uu^*=1\}.
\]
A self-adjoint unitary will be called a \emph{symmetry}.




\medskip
\noindent\textbf{Notation for families.}
If \(N\subset M\) is a von Neumann subalgebra and \(F\subset M\), the notation
\(E_N(F)=0\) means that \(E_N(f)=0\) for every \(f\in F\).  If
\(F,G\subset L^2(M,\tau)\), or if one of them is a closed subspace, then
\(F\perp G\) means that every element of \(F\) is orthogonal to every element
of \(G\).  We write \(F\dot\cup G\) for a disjoint union and \(\mathcal P(F)\)
for the power set of \(F\).  If \(\mathcal P_0\subset M\), then
\(\operatorname{*-alg}_{\mathbb Q+i\mathbb Q}(\mathcal P_0)\) denotes the
unital \(*\)-algebra over \(\mathbb Q+i\mathbb Q\) generated by
\(\mathcal P_0\).







In 1967, at the Baton Rouge conference~\cite{Kad67}, Kadison asked the
trace-vector orthonormal basis problem for \(\mathrm{II}_1\) factors in the
context of the natural Hilbert-space representations of group and
group-measure-space factors\footnote{Ge records Kadison's Problem~13 and
states that the two preceding group/GMS realization questions have negative
answers while the trace-vector basis question remains open; see
\cite[p.~622, Problem~13]{Ge03}.  See also \cite[Problem~S.2]{Pet20}.}.
We formulate the question in the tracial standard representation:

\begin{prob}[Kadison, tracial standard form]\label{prob:kadison}
Let \(M\) be a type \(\mathrm{II}_1\) factor with faithful normal tracial
state \(\tau\). Does \(L^2(M,\tau)\), for the left action of \(M\), admit a
Hilbert orthonormal basis consisting of trace vectors?
\end{prob}

In the tracial standard representation, Problem~\ref{prob:kadison} is also naturally reformulated as a
\emph{unitary basis problem}.  Indeed, if \(u\in\mathcal U(M)\), then
\[
   \langle x\widehat u,\widehat u\rangle_2
   =
   \tau(u^*xu)
   =
   \tau(x),\qquad x\in M,
\]
so \(\widehat u\) is a trace vector.  Conversely, as recalled in
Proposition~\ref{prop:trace-vectors-unitaries}, every trace vector in
\(L^2(M,\tau)\) is represented by a unitary in \(M\).  Thus,
Problem~\ref{prob:kadison} is equivalent to asking whether
\(L^2(M,\tau)\) has a Hilbert orthonormal basis consisting of unitaries in \(M\).

    
The separable case of this problem was settled in \cite{HTZ26}: using a
finite-dimensional norming lemma derived from the noncommutative Lyapunov
theorem of Akemann and Weaver~\cite{AW03}, it was shown that every diffuse finite von
Neumann algebra with separable $L^2$-space admits a Hilbert orthonormal basis
consisting of symmetries.

The purpose of the present paper is to prove the complementary nonseparable
case.  This case requires a different mechanism.  One must extend bases
through subalgebras of smaller density character while controlling
orthogonality to previously chosen, possibly very large, backgrounds.  The
finite-layer certified background formalism records enough von Neumann
algebraic structure to ensure that the relevant Hilbert-space projections
remain bounded elements of $M$.  This allows a transfinite extension argument
over the density character of $L^2(M,\tau)$.



\medskip
\noindent\textbf{Organization.}
Section~\ref{sec:preliminaries} contains the reduction from trace vectors to
unitaries and from trace-zero symmetry bases to full complex bases.
Sections~\ref{sec:relative-norming}--\ref{sec:all-cardinal-extension}
develop the relative extension machinery needed in the nonseparable setting.
Section~\ref{sec:nonseparable-main} proves the nonseparable theorem.
Section~\ref{sec:consequence} records the consequence obtained by combining
the present theorem with the separable result of \cite{HTZ26}.

\section{Preliminaries and reductions}\label{sec:preliminaries}

For a finite von Neumann algebra $(P,\tau)$, we shall write
\[
   H_0(P)
   =
   L^2(P,\tau)_{\sa}\ominus \RR\one
   =
   \{\xi\in L^2(P,\tau)_{\sa}:\langle \xi,\one\rangle_2=0\}.
\]
Thus, if $x=x^*\in P$, then $\widehat{x}\in H_0(P)$ exactly when
$\tau(x)=0$.  We also write
\[
   \mathcal S_0(P)
   =
   \{s\in P:s=s^*=s^{-1},\ \tau(s)=0\}
\]
for the set of trace-zero symmetries in $P$.  When $P$ is a von Neumann
subalgebra of a fixed finite von Neumann algebra, these notations are always
understood with respect to the restricted trace.

We shall use the standard finite von Neumann algebra facts that the tracial
standard representation has commutant given by right multiplication, that a
trace-preserving conditional expectation exists for every unital inclusion of
finite von Neumann algebras and is the \(L^2\)-orthogonal projection, and that
Kaplansky's density theorem gives \(L^2\)-density on bounded sets in the
presence of a finite trace; see, for example, \cite{Tak72,Tak02,KR97I,KR97II}.

We first record the standard reduction from trace vectors to unitaries.

\begin{prop}\label{prop:trace-vectors-unitaries}
Let \((M,\tau)\) be a finite von Neumann algebra in its tracial standard
representation on \(L^2(M,\tau)\). A vector \(\xi\in L^2(M,\tau)\) is a trace
vector for \(M\) if and only if
\[
   \xi=\widehat{u}
\]
for some unitary \(u\in\mathcal U(M)\).
\end{prop}
	
	\begin{proof}
	If \(u\in\mathcal U(M)\), then for every \(x\in M\),
	\[
	   \langle x\widehat{u},\widehat{u}\rangle_2
	   =\tau(u^*xu)
	   =\tau(x),
	\]
	so \(\widehat{u}\) is a trace vector.
	
Conversely, suppose that \(\xi\) is a trace vector. Let
\(\Omega=\widehat{1}\). Define
\[
   V_0:M\Omega\to L^2(M,\tau),
   \qquad
   V_0(x\Omega)=x\xi .
\]
For \(x\in M\),
	\[
	   \|V_0(x\Omega)\|_2^2
	   =\|x\xi\|_2^2
	   =\langle x^*x\xi,\xi\rangle_2
	   =\tau(x^*x)
	   =\|x\Omega\|_2^2.
	\]
	Thus \(V_0\) extends to an isometry \(V\) on \(L^2(M,\tau)\). Moreover, for
	\(a,x\in M\),
	\[
	   V(ax\Omega)=ax\xi=aV(x\Omega),
	\]
	so \(V\in M'\). Recall the standard identification of the commutant in the tracial
representation~\cite{Tak02}: in the tracial standard representation of
\(M\) on \(L^2(M,\tau)\), the commutant \(M'\) is the right regular
representation of \(M\): for each \(a\in M\),
\[
   R_a\widehat{x}=\widehat{xa},\qquad x\in M,
\]
and the map \(a\mapsto R_a\) is a \(*\)-anti-isomorphism from \(M\) onto
\(M'\). In particular \(M'\cong M^{\mathrm{op}}\) is finite.  Hence every isometry in
\(M'\) is unitary.  Indeed, let \(\tau'\) be the faithful normal finite trace
on \(M'\) transported from \(M^{\mathrm{op}}\).  If \(v\in M'\) and
\(v^*v=\one\), then
\[
   \tau'(\one-vv^*)
   =
   \tau'(\one)-\tau'(vv^*)
   =
   \tau'(v^*v)-\tau'(vv^*)
   =
   0 .
\]
Since \(\one-vv^*\ge0\) and \(\tau'\) is faithful, we get \(vv^*=\one\).
Thus \(V\) is a unitary element of \(M'\). Identifying
	\(M'\) with the right regular representation of \(M\), we get
	\(V\Omega=\widehat{u}\) for some \(u\in\mathcal U(M)\). But \(V\Omega=\xi\).
	Therefore \(\xi=\widehat{u}\).
	\end{proof}

The next lemma is the basic trace-zero reduction used in the nonseparable proof.

\begin{lem}[Reduction to trace-zero symmetries]\label{lem:trace-zero-reduction}
Let $M$ be a $\mathrm{II}_1$ factor.  If $B\subset\mathcal S_0(M)$ is a Hilbert orthonormal basis of the real Hilbert space $H_0(M)$, then $\{\widehat{\one}\}\cup\{\widehat{s}:s\in B\}$ is a Hilbert orthonormal basis of the complex Hilbert space $L^2(M,\tau)$.  Moreover every vector in $\{\widehat{\one}\}\cup\{\widehat{s}:s\in B\}$ is a trace vector for the left action of $M$.
\end{lem}

\begin{proof}
For self-adjoint $a,b\in M$, the number $\tau(ba)$ is real, since
\[
   \tau(ba)=\tau(ab)=\tau((ba)^*).
\]
Thus real orthogonality inside $L^2(M,\tau)_{\sa}$ is the same as complex
Hilbert-space orthogonality for these vectors.  Since
\(B\) is an orthonormal basis of \(H_0(M)\) and \(B\subset\mathcal S_0(M)\),
the family \(\{\one\}\cup B\) is orthonormal in \(L^2(M,\tau)\).

Let \(\mathcal L\) be the closed complex linear span of \(\{\one\}\cup B\).
Since
\[
   L^2(M,\tau)_{\sa}
   =
   \RR\one\oplus H_0(M)
   =
   \overline{\spanop_{\RR}}\bigl(\{\one\}\cup B\bigr),
\]
we have \(L^2(M,\tau)_{\sa}\subset \mathcal L\).  Hence
\(M_{\sa}\subset\mathcal L\), and therefore
\[
   M=M_{\sa}+iM_{\sa}\subset\mathcal L .
\]
Since \(M\) is dense in \(L^2(M,\tau)\), it follows that
\(\mathcal L=L^2(M,\tau)\).  Thus \(\{\one\}\cup B\) is an orthonormal
basis of the complex Hilbert space \(L^2(M,\tau)\).

If $s=s^*=s^{-1}\in M$, then for every $x\in M$,
\[
   \langle xs,s\rangle_2=\tau(sxs)=\tau(xss)=\tau(x).
\]
Thus every self-adjoint unitary is a trace vector, and so is $\one$.
\end{proof}



\section{A relative norming lemma under small-density constraints}
\label{sec:relative-norming}

Throughout Sections~\ref{sec:relative-norming}--\ref{sec:nonseparable-main},
$M$ denotes a $\mathrm{II}_1$ factor with faithful normal tracial state
$\tau$ and
\[
   \kappa=\dens L^2(M,\tau)>\aleph_0.
\]
All von Neumann subalgebras are understood to be unital, except that corners such as $eMe$ are taken with their natural unit $e$.  We use density character in the Hilbert norm: $\dens L^2(M)$ is the least cardinality of an $L^2$-dense subset.

We shall use standard facts on finite von Neumann algebras without further
comment: Kaplansky's density theorem, the trace-preserving normal conditional
expectation onto a von Neumann subalgebra of a finite von Neumann algebra, and
projection comparison in finite factors; see, for example,
\cite{Tak72,KR97II,Tak02}.

\begin{lem}[Elementary $L^2$-closure and cardinal facts]\label{lem:l2-cardinal-facts}
Let $(P,\tau)$ be a finite von Neumann algebra.

\begin{enumerate}[label=(\roman*)]
\item Let $A_0\subset P$ be a unital $C^*$-subalgebra and let
      $A=A_0''=W^*(A_0)$.  If $A_0$ has norm density at most $\mu$, then
      $L^2(A,\tau)$ has density at most $\max\{\mu,\aleph_0\}$.

\item If $(P_i)_{i\in I}$ is a directed increasing family of von Neumann
      subalgebras of $P$ and
      \[
         P_0=W^*\Bigl(\bigcup_{i\in I}P_i\Bigr),
      \]
      then
      \[
         L^2(P_0,\tau)=
         \overline{\bigcup_{i\in I}L^2(P_i,\tau)}^{\|\cdot\|_2}.
      \]

\item If $Q\subset P$ is a von Neumann subalgebra and $x\in P$ satisfies
      $\widehat{x}\in L^2(Q,\tau)$, then $x\in Q$.

\item If $Y\subset P_{\sa}$ is $L^2$-dense in $L^2(P,\tau)_{\sa}$, then
      $W^*(Y)=P$.
\end{enumerate}
\end{lem}

\begin{proof}
For (i), choose a norm-dense subset of $A_0$ of cardinal at most
$\max\{\mu,\aleph_0\}$ and close it under rational complex linear
combinations, adjoints, and products.  This gives a norm-dense $*$-subalgebra
of $A_0$ of cardinal at most $\max\{\mu,\aleph_0\}$.  By Kaplansky's density
theorem, the unit ball of $A_0$ is strong-$*$ dense in the unit ball of $A$.
On bounded subsets of a finite von Neumann algebra, strong convergence in the
standard representation implies $L^2$-convergence, because
\[
   \|x_\alpha-x\|_2=\|(x_\alpha-x)\widehat{1}\|_2 .
\]
Hence the unit ball of $A$ has an $L^2$-dense subset of cardinal at most
$\max\{\mu,\aleph_0\}$, and scaling gives the assertion for $L^2(A)$.

For (ii), put
\[
   B=\overline{\bigcup_iP_i}^{\|\cdot\|}.
\]
Then \(B\) is a \(C^*\)-algebra and \(B''=P_0\).  By Kaplansky's density
theorem, the unit ball of \(B\) is strong-\(*\) dense in the unit ball of
\(P_0\).  Since \(\bigcup_iP_i\) is norm dense in \(B\), it follows that
\(\bigcup_iP_i\) is \(L^2\)-dense in \(P_0\) on bounded sets, and hence
\[
   L^2(P_0,\tau)=
   \overline{\bigcup_iL^2(P_i,\tau)}^{\|\cdot\|_2}.
\]

For (iii), let $E_Q:P\to Q$ be the trace-preserving conditional expectation.
Since $E_Q$ is the $L^2$-orthogonal projection onto $L^2(Q)$, the assumption
implies $\widehat{x}=\widehat{E_Q(x)}$.  Faithfulness of $\tau$ gives
$x=E_Q(x)\in Q$.

For (iv), put $Q=W^*(Y)$.  Then $L^2(Q)_{\sa}$ is a closed real subspace of
$L^2(P)_{\sa}$ containing $Y$, hence $L^2(Q)_{\sa}=L^2(P)_{\sa}$.  Therefore
$L^2(Q)=L^2(P)$, and (iii) gives $Q=P$.
\end{proof}

\begin{lem}[Density of nonzero corners]\label{lem:corner-density}
Let $M$ be a $\mathrm{II}_1$ factor with $\dens L^2(M)=\kappa>\aleph_0$.  If $0\neq e\in M$ is a projection, then
\[
   \dens L^2(eMe)=\kappa,
\]
where $eMe$ is equipped with its finite trace inherited from $\tau$.
\end{lem}

\begin{proof}
Clearly $\dens L^2(eMe)\leq\kappa$.  Let $\mu=\dens L^2(eMe)$.  Choose $n\in\NN$ with $1/n\leq\tau(e)$.  In a $\mathrm{II}_1$ factor there are orthogonal projections $e_1,\ldots,e_n$ with $\sum_i e_i=\one$ and $\tau(e_i)=1/n$.  There are also projections $f_i\leq e$ with $\tau(f_i)=1/n$; these $f_i$ need not be mutually orthogonal.  Let $v_i$ be partial isometries with $v_i^*v_i=f_i$ and $v_iv_i^*=e_i$.  For $x\in M$,
\[
   e_i x e_j=v_i\,(v_i^*xv_j)\,v_j^*,
   \qquad v_i^*xv_j\in f_iMf_j\subset eMe.
\]
Thus $M$ is contained in the finite linear span of the subspaces $v_i(eMe)v_j^*$.  Each of these subspaces has $L^2$-density at most $\mu$, so
\[
   \kappa=\dens L^2(M)\leq\max\{\mu,\aleph_0\}.
\]
Since $\kappa>\aleph_0$, this inequality forces $\mu>\aleph_0$ and hence $\max\{\mu,\aleph_0\}=\mu$.  Therefore $\kappa\leq\mu$, and the reverse inequality already observed gives $\mu=\kappa$.
\end{proof}

\begin{lem}[Small-constraint perturbation]\label{lem:perturb}
Let $\Gamma\subset M_{\sa}$ satisfy $|\Gamma|<\kappa$.  If $0\neq e\in M$ is a projection, then there exists a nonzero $h=h^*\in eMe$ such that
\[
   \tau(\gamma h)=0,
   \qquad \gamma\in\Gamma.
\]
\end{lem}

\begin{proof}
Let $D=W^*(\{e\}\cup\Gamma)$.  If $\lambda=|\Gamma|$, then the unital
$C^*$-algebra generated by $\{e\}\cup\Gamma$ has norm density at most
$\max\{\lambda,\aleph_0\}<\kappa$, because rational-coefficient
$*$-polynomials in the generators are norm dense in it.  By
Lemma~\ref{lem:l2-cardinal-facts}(i),
\[
   \dens L^2(D)<\kappa.
\]
Hence also $\dens L^2(eDe)<\kappa$.

By Lemma~\ref{lem:corner-density},
\[
   \dens L^2(eMe)=\kappa.
\]
Since
\[
   L^2(eMe)=L^2(eMe)_{\sa}+iL^2(eMe)_{\sa}
\]
and similarly for \(eDe\), the real Hilbert space
\(L^2(eMe)_{\sa}\) has density \(\kappa\), while
\(L^2(eDe)_{\sa}\) has density \(<\kappa\).  Thus
\(L^2(eDe)_{\sa}\) is a proper closed real subspace of
\(L^2(eMe)_{\sa}\). Choose
\[
   \eta=\eta^*\in L^2(eMe)_{\sa}\setminus L^2(eDe)_{\sa}.
\]
Since $(eMe)_{\sa}$ is $L^2$-dense in $L^2(eMe)_{\sa}$ and
$L^2(eDe)_{\sa}$ is closed, we may choose
\(a=a^*\in eMe\) so close to \(\eta\) in \(L^2\)-norm that
\[
   \widehat a\notin L^2(eDe)_{\sa}.
\]
In particular,
\[
   a\in eMe\setminus eDe .
\]
Let
\[
   h=a-E_D(a),
\]
where $E_D:M\to D$ is the trace-preserving conditional expectation. Since $e\in D$ and $a=eae$, bimodularity of $E_D$ gives
\[
   E_D(a)=E_D(eae)=eE_D(a)e\in eDe.
\]
Thus $h\in eMe$.  Moreover $h\neq0$, because otherwise $a=E_D(a)\in eDe$.

Finally, for $\gamma\in\Gamma\subset D$, trace preservation and
$D$-bimodularity of $E_D$ give
\[
   \tau(\gamma a)
   =\tau(E_D(\gamma a))
   =\tau(\gamma E_D(a)).
\]
Therefore $\tau(\gamma h)=0$ for every $\gamma\in\Gamma$.
\end{proof}

\begin{lem}[Large relative norming lemma]\label{lem:large-relative-norming}
Let $\kappa=\dens L^2(M)>\aleph_0$.  Let $N\subset M$ be a von Neumann subalgebra with $\dens L^2(N)<\kappa$, and let $\Gamma\subset M_{\sa}$ have cardinality $<\kappa$.  Suppose $0\neq r=r^*\in M$ satisfies
\[
   E_N(r)=0,
   \qquad
   \tau(\gamma r)=0\quad(\gamma\in\Gamma).
\]
Then there exists $s\in\mathcal S_0(M)$ such that
\[
   E_N(s)=0,
   \qquad
   \tau(\gamma s)=0\quad(\gamma\in\Gamma),
\]
and
\[
   \langle r,s\rangle_2=\frac{\|r\|_2^2}{\|r\|_\infty}.
\]
\end{lem}

\begin{proof}
Set $c=\|r\|_\infty$. Since $r\neq0$, we have $c>0$.  Put
\[
   q=\frac{\one+r/c}{2}.
\]
Then $0\leq q\leq\one$.  Choose a subset $\mathcal D\subset N_{\sa}$ consisting of bounded operators, $L^2$-dense in $L^2(N)_{\sa}$, and satisfying
\[
   |\mathcal D|\leq\max\{\dens L^2(N),\aleph_0\}<\kappa.
\]
Consider the set
\[
K=\Bigl\{0\leq x=x^*\leq\one:\
\begin{array}{l}
\tau(yx)=\tau(yq)\quad (y\in\mathcal D),\\
\tau(\gamma x)=\tau(\gamma q)\quad (\gamma\in\Gamma),\\
\tau(rx)=\tau(rq)
\end{array}
\Bigr\}.
\]
It is nonempty because $q\in K$.  It is ultraweakly compact and convex: indeed, it is the intersection of the ultraweakly compact positive unit ball of $M$ with ultraweakly closed affine hyperplanes defined by normal functionals.  By the Krein--Milman theorem~\cite{Con90}, $K$ has an extreme point.  Let $p$ be one.

We claim that $p$ is a projection.  Suppose not.  Since $0\leq p\leq1$ and
$p$ is not a projection, the spectral projection of $p$ corresponding to the
open interval $(0,1)$ is nonzero.  Since
\[
   (0,1)=\bigcup_{n=2}^\infty [1/n,1-1/n],
\]
there is some $0<\varepsilon<1/2$ such that
\[
   e=1_{[\varepsilon,1-\varepsilon]}(p)
\]
is nonzero.

Apply Lemma~\ref{lem:perturb} to the set
\[
   \mathcal D\cup\Gamma\cup\{r\}
\]
and to the projection $e$.  This set has cardinal $<\kappa$.  We get a
nonzero $h=h^*\in eMe$ such that
\[
   \tau(yh)=0\quad(y\in\mathcal D),\qquad
   \tau(\gamma h)=0\quad(\gamma\in\Gamma),\qquad
   \tau(rh)=0.
\]
After replacing $h$ by a sufficiently small nonzero real scalar multiple, we
may assume $\|h\|_\infty<\varepsilon$.  Since \(e\) is a spectral projection
of \(p\), it commutes with \(p\).  Also \(h=ehe\).  Hence, with respect to
the decomposition \(eL^2(M)\oplus(1-e)L^2(M)\), the operators \(p\pm h\)
are block diagonal:
\[
   p\pm h=(pe\pm h)\oplus p(1-e).
\]
On \(eL^2(M)\) we have
\[
   \varepsilon e\leq pe\leq(1-\varepsilon)e.
\]
Since \(h=h^*\in eMe\), we have
\[
   -\|h\|_\infty e\le h\le \|h\|_\infty e.
\]
As \(\|h\|_\infty<\varepsilon\), it follows that
\[
   0\leq pe\pm h\leq e.
\]
On \((1-e)L^2(M)\), the operator \(p(1-e)\) is a positive contraction.
Therefore
\[
   0\leq p\pm h\leq1.
\]
The displayed orthogonality relations show that $p+h$ and $p-h$ satisfy all
defining affine constraints of $K$.  Thus $p\pm h\in K$, and
\[
   p=\frac12(p+h)+\frac12(p-h)
\]
is a nontrivial convex decomposition in $K$, contradicting the extremality
of $p$.  Hence $p$ is a projection.

Put $s=2p-\one$.  For $y\in\mathcal D$,
\[
   \tau(yq)=\frac12\tau(y)+\frac{1}{2c}\tau(yr)=\frac12\tau(y),
\]
since $E_N(r)=0$. Thus $\tau(ys)=0$ for every $y\in\mathcal D$.  If
$z\in N_{\sa}$, choose $y_i\in\mathcal D$ with
$\|y_i-z\|_2\to0$.  Since $s\in M$ is bounded, the functional
$w\mapsto\tau(ws)$ is $L^2$-continuous on $L^2(N)_{\sa}$.  Hence
$\tau(zs)=0$ for all $z\in N_{\sa}$. Since \(E_N(s)\in N_{\sa}\), we may take \(z=E_N(s)\).  This gives
\[
0=\tau(E_N(s)s)
  =\tau(E_N(E_N(s)s))
  =\tau(E_N(s)^2)
  =\|E_N(s)\|_2^2,
\]
where we used that \(E_N\) is trace preserving and \(N\)-bimodular.  Hence
\(E_N(s)=0\).  Similarly, for \(\gamma\in\Gamma\),
\[
   \tau(\gamma q)=\frac12\tau(\gamma)+\frac{1}{2c}\tau(\gamma r)
   =\frac12\tau(\gamma),
\]
so \(\tau(\gamma s)=0\).  Since \(\one\in N\), the identity \(E_N(s)=0\)
also gives \(\tau(s)=0\).  Therefore \(s\in\mathcal S_0(M)\).

Finally, using the convention that the inner product is linear in the first
variable and using traciality,
\[
\begin{aligned}
   \langle r,s\rangle_2
   &=\tau(sr)
     =\tau(rs) \\
   &=\tau(r(2p-\one)) \\
   &=2\tau(rp)-\tau(r) \\
   &=2\tau(rq)-\tau(r) \\
   &=2\left(\frac12\tau(r)+\frac{1}{2c}\tau(r^2)\right)-\tau(r)
     =\frac{\|r\|_2^2}{\|r\|_\infty}.
\end{aligned}
\]
\end{proof}

\section{Completed blocks and certified backgrounds}\label{sec:certified-backgrounds}

\begin{defin}[Completed block]\label{def:block}
A \emph{completed block} in $M$ is a triple
\[
   \mathfrak B=(A,Q,T),
\]
where $A\subset Q\subset M$ are von Neumann subalgebras and
$T\subset\mathcal S_0(M)$ is an orthonormal family such that
\[
   L^2(Q)_{\sa}
   =
   L^2(A)_{\sa}\oplus\overline{\spanop_{\RR}}T .
\]
The sum is an orthogonal direct sum of real Hilbert spaces.

If a von Neumann subalgebra $N\subset M$ is specified, we say that
$\mathfrak B$ is a \emph{completed block over $N$} if, in addition,
\[
   A\subset N
   \qquad\text{and}\qquad
   E_N(t)=0\quad(t\in T).
\]
In that case $T\perp L^2(N)_{\sa}$, and in particular
$T\perp L^2(A)_{\sa}$.

In either case, the displayed decomposition implies $T\subset Q$: indeed
$t\in T$ is an element of $M$ whose $L^2$-vector belongs to $L^2(Q)$, so
Lemma~\ref{lem:l2-cardinal-facts}(iii) gives $t\in Q$.
\end{defin}

\begin{lem}[Projection formula]\label{lem:block-projection}
If $\mathfrak B=(A,Q,T)$ is a completed block, then for every $x=x^*\in M$ the orthogonal projection of $x$ onto $\overline{\spanop_{\RR}}T$ is
\[
   P_T^{\mathfrak B}(x)=E_Q(x)-E_A(E_Q(x)).
\]
In particular $P_T^{\mathfrak B}(x)\in M_{\sa}$.
\end{lem}

\begin{proof}
The map $E_Q$ is the $L^2$-orthogonal projection onto $L^2(Q)$ and $E_A$ is the $L^2$-orthogonal projection from $L^2(Q)$ onto $L^2(A)$.  The completed-block decomposition identifies $L^2(Q)_{\sa}\ominus L^2(A)_{\sa}$ with $\overline{\spanop_{\RR}}T$.
\end{proof}

We use finite certificates to record when projections onto large orthonormal families are represented by bounded operators.  The word ``finite'' refers to the depth of the certificate tree; the individual layers may be uncountable.

\begin{defin}[Finite-layer certificates and certified backgrounds]\label{def:certified}
A \emph{finite-layer certificate} $\mathfrak c$ is defined recursively,
together with its underlying orthonormal family
\[
   \Omega(\mathfrak c)\subset\mathcal S_0(M)
\]
and, for every $x=x^*\in M$, a bounded self-adjoint element
$P_{\mathfrak c}(x)\in M_{\sa}$.

\smallskip

\noindent\emph{The empty certificate.}
There is an empty certificate $\varnothing$ with
\[
   \Omega(\varnothing)=\varnothing,
   \qquad
   P_{\varnothing}(x)=0 .
\]

\smallskip

\noindent\emph{Adding one layer.}
Suppose that $\mathfrak c^-$ is a finite-layer certificate with
\[
   \Omega^-=\Omega(\mathfrak c^-).
\]
Suppose also that $\mathfrak c_\Theta$ is a finite-layer certificate with
\[
   \Theta=\Omega(\mathfrak c_\Theta)\subset\Omega^- .
\]
Let $U\subset\mathcal S_0(M)$ be an orthonormal family which is orthogonal to
$\Omega^-$.  Assume that there are von Neumann subalgebras
\[
   C\subset R\subset M
\]
such that
\[
   L^2(R)_{\sa}
   =
   L^2(C)_{\sa}\oplus
   \overline{\spanop_{\RR}}(\Theta\cup U).
\]
Then the finite data
\[
   \mathfrak c=
   \bigl(\mathfrak c^-,\mathfrak c_\Theta,C,R,U\bigr)
\]
is a finite-layer certificate with underlying family
\[
   \Omega(\mathfrak c)=\Omega^-\dot\cup U.
\]
We say that $U$ is the \emph{new layer} of $\mathfrak c$, and that it is a
layer over the certified subbackground $(\Theta,\mathfrak c_\Theta)$.

For $x=x^*\in M$, define
\[
   P_U^{\mathfrak c}(x)
   =
   \bigl(E_R(x)-E_C(E_R(x))\bigr)-P_{\mathfrak c_\Theta}(x),
\]
and
\[
   P_{\mathfrak c}(x)
   =
   P_{\mathfrak c^-}(x)+P_U^{\mathfrak c}(x).
\]

A \emph{finite-layer certified background} is a pair
$(\Omega,\mathfrak c)$, where $\mathfrak c$ is a finite-layer certificate and
$\Omega=\Omega(\mathfrak c)$.  When the certificate is fixed, we suppress it
from the notation and write simply $P_\Omega(x)$ for $P_{\mathfrak c}(x)$.

A subfamily $\Theta\subset\Omega$ is called a \emph{certified subbackground}
when it is equipped with some finite-layer certificate
$\mathfrak c_\Theta$ satisfying $\Omega(\mathfrak c_\Theta)=\Theta$.
This certificate need not be a literal subcertificate of the fixed certificate
for $\Omega$; it is part of the finite data being used.  The adjective
\emph{finite-layer} refers only to the finite depth of the certificate.  The
individual layers may have arbitrary cardinality.
\end{defin}

\begin{rem}[Complexity of certificates]\label{rem:certificate-complexity}
We define the complexity of a finite-layer certificate recursively by
\[
   \ell(\varnothing)=0,
\]
and, for
\[
   \mathfrak c=(\mathfrak c^-,\mathfrak c_\Theta,C,R,U),
\]
by
\[
   \ell(\mathfrak c)
   =
   1+\ell(\mathfrak c^-)+\ell(\mathfrak c_\Theta).
\]
All inductions on finite-layer certificates below are taken with respect to
this finite complexity.
\end{rem}

\begin{lem}[Boundedness of certified projections]\label{lem:certified-projection}
Let $\Omega$ be a finite-layer certified background.  Then for every $x=x^*\in M$, the Hilbert-space projection $P_\Omega x$ onto $\overline{\spanop_{\RR}}\Omega$ belongs to $M_{\sa}$.  More precisely, $P_\Omega x$ is a finite sum of expressions obtained from $x$ by trace-preserving conditional expectations appearing in the certificate.
\end{lem}

\begin{proof}
We prove the assertion by induction on the complexity of the finite certificate.
The empty certificate is trivial.

Let
\[
   \mathfrak c=
   \bigl(\mathfrak c^-,\mathfrak c_\Theta,C,R,U\bigr)
\]
be a certificate obtained by adding one layer.  Put
\[
   \Omega^-=\Omega(\mathfrak c^-),
   \qquad
   \Theta=\Omega(\mathfrak c_\Theta),
   \qquad
   \Omega=\Omega^-\dot\cup U.
\]
By the induction hypothesis, $P_{\mathfrak c^-}(x)$ is the Hilbert-space
projection of $x$ onto $\overline{\spanop_{\RR}}\Omega^-$, and
$P_{\mathfrak c_\Theta}(x)$ is the Hilbert-space projection of $x$ onto
$\overline{\spanop_{\RR}}\Theta$.  Both belong to $M_{\sa}$.

The completed-block decomposition
\[
   L^2(R)_{\sa}
   =
   L^2(C)_{\sa}\oplus
   \overline{\spanop_{\RR}}(\Theta\cup U)
\]
and Lemma~\ref{lem:block-projection} show that
\[
   E_R(x)-E_C(E_R(x))
\]
is the Hilbert-space projection of $x$ onto
$\overline{\spanop_{\RR}}(\Theta\cup U)$.  Since $U\perp\Omega^-$ and
$\Theta\subset\Omega^-$, this projection splits orthogonally as the sum of
the projection onto $\overline{\spanop_{\RR}}\Theta$ and the projection onto
$\overline{\spanop_{\RR}}U$.  Hence
\[
   P_U^{\mathfrak c}(x)
   =
   \bigl(E_R(x)-E_C(E_R(x))\bigr)-P_{\mathfrak c_\Theta}(x)
\]
is exactly the Hilbert-space projection of $x$ onto
$\overline{\spanop_{\RR}}U$.  Therefore
\[
   P_{\mathfrak c}(x)
   =
   P_{\mathfrak c^-}(x)+P_U^{\mathfrak c}(x)
\]
is the Hilbert-space projection of $x$ onto
$\overline{\spanop_{\RR}}\Omega$.  The displayed formula is a finite
combination of trace-preserving conditional expectations and previously
constructed projection formulas, so it belongs to $M_{\sa}$.
\end{proof}


Consequently, if two finite-layer certificates have the same underlying
orthonormal family \(\Theta\), then their formulas give the same element of
\(M_{\sa}\) for every \(x=x^*\in M\), since both compute the Hilbert-space
projection of \(x\) onto \(\overline{\spanop_{\RR}}\Theta\).


\section{The strong closure lemma for certified backgrounds}\label{sec:strong-closure}

The main technical issue is that an arbitrary subfamily of a certified background need not have bounded projection.  The next proposition gives a canonical way to enlarge any small subfamily to a small certified subbackground.



\medskip
\noindent\textbf{Purpose of the closure construction.}
The closure operation below is used only to ensure the following point: whenever
a small subfamily of a certified background is selected, it can be enlarged to
a still small subfamily on which all relevant projection supports are closed.
Consequently the Hilbert-space projections onto these enlarged subfamilies are
represented by bounded elements of the ambient von Neumann algebra \(M\).  No
additional model-theoretic property of \(M\) is being assumed.



For an orthonormal family $T$ and $\xi\in L^2(M)$, write
\[
   \supp_T(\xi)=\{t\in T:\langle \xi,t\rangle_2\neq0\}.
\]

\begin{lem}[Countability of supports]\label{lem:countable-supports}
Let $T$ be an orthonormal family in a Hilbert space $H$.  For every
$\xi\in H$,
\[
   \supp_T(\xi)=\{t\in T:\langle \xi,t\rangle\neq0\}
\]
is countable.
\end{lem}

\begin{proof}
For $n\geq1$, the set
\[
   F_n=\{t\in T:|\langle \xi,t\rangle|\geq 1/n\}
\]
is finite by Bessel's inequality.  Since
\[
   \supp_T(\xi)=\bigcup_{n=1}^\infty F_n,
\]
the support is countable.
\end{proof}


Here and below, for a regular cardinal \(\chi\), \(H(\chi)\) denotes the set
of all sets whose transitive closure has cardinality \(<\chi\).  We use only
the standard elementary consequences of closing \(H(\chi)\) under countably
many finitary Skolem functions.

\begin{lem}[Skolem-hull closure fact]\label{lem:skolem-hull}
Let \(\chi\) be regular and let \(\mathfrak A\) be a countable first-order
expansion of \((H(\chi),\in)\) by finitary function symbols, including
Skolem functions, constants for \(\omega\), and constants for all natural
numbers.  For \(S\subset H(\chi)\), let \(\Hull_{\mathfrak A}(S)\) denote the
Skolem hull generated by \(S\) together with the named constants.  Then:
\begin{enumerate}[label=(\roman*)]
\item \(|\Hull_{\mathfrak A}(S)|\le |S|+\aleph_0\);
\item \(S_1\subset S_2\) implies
      \(\Hull_{\mathfrak A}(S_1)\subset\Hull_{\mathfrak A}(S_2)\);
\item for every increasing chain \((S_i)_{i\in I}\),
      \[
      \Hull_{\mathfrak A}\Bigl(\bigcup_i S_i\Bigr)=
      \bigcup_i \Hull_{\mathfrak A}(S_i);
      \]
\item if \(c\in\Hull_{\mathfrak A}(S)\) is countable in the universe, then
      \(c\subset\Hull_{\mathfrak A}(S)\).
\end{enumerate}
\end{lem}

\begin{proof}
The first three assertions follow from the usual construction of the closure
under countably many finitary Skolem functions; in the chain-continuity
statement, every finite tuple from an increasing union is already contained in
one member of the chain. For (iv), the hull is an elementary submodel of the
Skolemized structure and contains every natural number.  If \(c=\varnothing\),
there is nothing to prove.  Assume that \(c\neq\varnothing\).  Since \(c\) is
countable, \(H(\chi)\) satisfies that there is a surjection
\(f:\omega\to c\).  By elementarity and the Skolem functions, such an \(f\)
belongs to the hull.  For every \(n\in\omega\), the value \(f(n)\) is definable
from \(f\) and \(n\), and hence belongs to the hull.  Thus
\(c\subset\Hull_{\mathfrak A}(S)\).  This is the standard Skolem-hull argument;
see, for example, \cite{Jech03,CK90}.
\end{proof}



\begin{lem}[Closed subblocks of a completed block]\label{lem:closed-subblock}
Let $\mathfrak B=(A,Q,T)$ be a completed block.  There is a closure operation $\cl_{\mathfrak B}$ on subsets of $T$ with the following properties:
\begin{enumerate}[label=(\roman*)]
\item $S\subset\cl_{\mathfrak B}(S)$;
\item $\cl_{\mathfrak B}(\cl_{\mathfrak B}(S))=\cl_{\mathfrak B}(S)$;
\item $|\cl_{\mathfrak B}(S)|\leq |S|+\aleph_0$;
\item $S_1\subset S_2$ implies $\cl_{\mathfrak B}(S_1)\subset\cl_{\mathfrak B}(S_2)$;
\item for every increasing chain $(S_i)_{i\in I}$,
\[
   \cl_{\mathfrak B}\Bigl(\bigcup_iS_i\Bigr)=\bigcup_i\cl_{\mathfrak B}(S_i);
\]
\item if $S=\cl_{\mathfrak B}(S)$, then there are von Neumann algebras $A_S\subset A$ and $Q_S\subset Q$ such that
\[
   (A_S,Q_S,S)
\]
is a completed block.
\end{enumerate}
\end{lem}

\begin{proof}
We spell out the Skolem-hull construction. Choose a regular cardinal \(\chi\) such that
\[
   M,A,Q,T,\tau,E_A,E_Q,\omega,\ [T]^{\leq\aleph_0}
\]
belong to \(H(\chi)\), and such that the graphs of the algebraic operations on
\(M\) also belong to \(H(\chi)\).  We take \(\chi\) large enough for all these
requirements throughout the argument below.

Work in a countable first-order expansion \(\mathfrak A_{\mathfrak B}\) of
\((H(\chi),\in)\).  The objects \(M,A,Q,T\), the algebraic operations on \(M\),
and the conditional expectations \(E_A,E_Q\) are named as single predicates,
function symbols, or set-valued parameters. The algebraic operations include the constants \(0_M\) and \(\one_M\),
addition, multiplication, adjoint, and scalar multiplication by each
\(\lambda\in\mathbb Q+i\mathbb Q\).  We do not add constants for all individual
elements of \(M,A,Q\), or \(T\).

We include a constant for \(\omega\) and constants for every natural number.
Hence every Skolem hull \(\mathcal E\) considered below satisfies
\[
   \omega\in\mathcal E
   \qquad\text{and}\qquad
   \omega\subset\mathcal E .
\]
Every partial map used below is made into a total function on \(H(\chi)\) by
assigning an arbitrary fixed value outside its intended domain.

As a function symbol in the expansion, include the total support map
\[
   \sigma_{\mathfrak B}(x)=
   \supp_T\bigl(E_Q(x)-E_A(E_Q(x))\bigr),
   \qquad x=x^*\in M,
\]
again with an arbitrary fixed value outside \(M_{\sa}\).  By
Lemma~\ref{lem:countable-supports}, for \(x=x^*\in M\) the value
\(\sigma_{\mathfrak B}(x)\) is a countable subset of \(T\), and hence belongs
to \(H(\chi)\).

The language is countable, since only countably many algebraic operations,
conditional expectations, support maps, and named structural parameters have
been added.  Add countably many finitary Skolem functions for this countable
structure.
For \(S\subset T\), let \(\Hull_{\mathfrak B}(S)\) be the Skolem hull generated
by \(S\), together with the fixed structural parameters just described, the
constant \(\omega\), and the constants for all natural numbers.  Define
\[
   \cl_{\mathfrak B}(S)=T\cap\Hull_{\mathfrak B}(S).
\]
Thus the hull is generated from \(S\), the natural numbers, and the fixed
structural objects, but not from all individual elements of \(M,A,Q\), or \(T\).

By Lemma~\ref{lem:skolem-hull}, the size, monotonicity, and chain-continuity
assertions follow from the countability of the language and the finitarity of
the Skolem functions.
Idempotence follows because
\[
   \Hull_{\mathfrak B}(S)
   \subset \Hull_{\mathfrak B}(\cl_{\mathfrak B}(S))
   \subset \Hull_{\mathfrak B}(\Hull_{\mathfrak B}(S))
   =\Hull_{\mathfrak B}(S).
\]

Assume now that $S=\cl_{\mathfrak B}(S)$, and set
$\mathcal E=\Hull_{\mathfrak B}(S)$.  Then $S=T\cap\mathcal E$.  We shall use
Lemma~\ref{lem:skolem-hull}\emph{(iv)}: if $c\in\mathcal E$ is countable, then
$c\subset\mathcal E$.

Define
\[
   A_S=W^*(A\cap\mathcal E),
   \qquad
   Q_S=W^*(A_S,S).
\]
We show that $(A_S,Q_S,S)$ is a completed block. Let \(w=w(a_1,\ldots,a_m,s_1,\ldots,s_n)\) be a bounded self-adjoint
\(*\)-word with coefficients in \(\mathbb Q+i\mathbb Q\), where
\(a_i\in A\cap\mathcal E\) and \(s_j\in S\). Since
\(S\subset T\subset Q\) and \(A\subset Q\), we have \(w\in Q\).  Moreover,
because the algebraic operations are included in the Skolemized structure,
\(w\in\mathcal E\).  In the original block,
\[
   w=E_A(w)+P_T^{\mathfrak B}(w).
\]
Since $E_A$ is part of the structure, $E_A(w)\in A\cap\mathcal E$. Moreover
\[
   \supp_T(P_T^{\mathfrak B}(w))=\sigma_{\mathfrak B}(w)\in\mathcal E .
\]
This set is countable.  Since \(\mathcal E\) contains \(\omega\) and all
natural numbers, Lemma~\ref{lem:skolem-hull}\emph{(iv)} gives
\[
   \sigma_{\mathfrak B}(w)\subset\mathcal E .
\]
Hence
\[
   \sigma_{\mathfrak B}(w)\subset T\cap\mathcal E=S,
\]
and therefore
\[
   P_T^{\mathfrak B}(w)\in\overline{\spanop_{\RR}}S.
\]
Let \(\mathcal A_S\) be the unital rational-coefficient algebraic
\(*\)-algebra generated by \(A\cap\mathcal E\) and \(S\), and let
\[
   B_S=\overline{\mathcal A_S}^{\|\cdot\|}
\]
be its norm closure.  Then
\[
   Q_S=W^*(A_S,S)=W^*(A\cap\mathcal E,S)=B_S'' .
\]
Kaplansky's density theorem applies to the $C^*$-algebra $B_S$: the
self-adjoint part of the unit ball of $B_S$ is strong-$*$ dense in the
self-adjoint part of the unit ball of $Q_S$.  Since $\mathcal A_S$ is norm
dense in $B_S$, and since bounded strong convergence implies
$L^2$-convergence in a finite von Neumann algebra, the bounded self-adjoint rational-coefficient \(*\)-words in
\(A\cap\mathcal E\) and \(S\) are \(L^2\)-dense in
\(L^2(Q_S)_{\sa}\). Hence
\[
   L^2(Q_S)_{\sa}\subset L^2(A_S)_{\sa}+\overline{\spanop_{\RR}}S.
\]
The reverse inclusion is clear from the definition of $Q_S$, and the sum is orthogonal because $S\perp L^2(A)_{\sa}$ and $A_S\subset A$.  Thus
\[
   L^2(Q_S)_{\sa}=L^2(A_S)_{\sa}\oplus\overline{\spanop_{\RR}}S.
\]
\end{proof}

\begin{prop}[Strong closure for finite-layer backgrounds]\label{prop:strong-closure}
Let $\Omega$ be a finite-layer certified background with a fixed finite certificate.  There is a closure operation
\[
   \cl_\Omega:\mathcal P(\Omega)\to\mathcal P(\Omega)
\]
such that:
\begin{enumerate}[label=(\roman*)]
\item $S\subset\cl_\Omega(S)$;
\item $\cl_\Omega(\cl_\Omega(S))=\cl_\Omega(S)$;
\item $|\cl_\Omega(S)|\leq |S|+\aleph_0$;
\item $S_1\subset S_2$ implies $\cl_\Omega(S_1)\subset\cl_\Omega(S_2)$;
\item for every increasing chain $(S_i)_{i\in I}$,
\[
   \cl_\Omega\Bigl(\bigcup_iS_i\Bigr)=\bigcup_i\cl_\Omega(S_i);
\]
\item if $S=\cl_\Omega(S)$, then $S$ is itself a finite-layer certified background.  Consequently $P_S(x)\in M_{\sa}$ for every $x=x^*\in M$.
\end{enumerate}
Moreover the closure operations may be chosen compatibly with adding one more layer: if $U$ is a new layer over a $\cl_\Omega$-closed certified subbackground $\Theta\subset\Omega$ and $\Omega^+=\Omega\dot\cup U$, then $\cl_{\Omega^+}$ can be chosen so that every $\cl_{\Omega^+}$-closed $\Xi\subset\Omega^+$ satisfies
\[
   \Xi\cap\Omega=\cl_\Omega(\Xi\cap\Omega).
\]
\end{prop}

\begin{proof}
We prove the closure assertions and the compatibility assertion simultaneously
by induction on the complexity of the fixed finite certificate. The empty case
is trivial. Suppose $\Omega=\Omega^-\dot\cup U$, where $\Omega^-$ has already been given a closure operation with the stated properties, and where $U$ is a layer over a certified subbackground $\Theta\subset\Omega^-$.  Thus there are von Neumann algebras $C\subset R\subset M$ with
\[
   L^2(R)_{\sa}=L^2(C)_{\sa}\oplus\overline{\spanop_{\RR}}(\Theta\cup U).
\]

Choose a large \(H(\chi)\) and a countable Skolemized language, with constants
for all natural numbers.  As in Lemma~\ref{lem:closed-subblock}, all
von Neumann algebras, orthonormal families, conditional expectations, support
maps, and previously chosen Skolem functions are named as single predicates, function symbols, or set-valued parameters; we do not name all their individual elements as constants.  The language contains:
\begin{itemize}[leftmargin=2em]
\item all data and Skolem functions used to define the already constructed closure operations for $\Omega^-$ and for every certified subbackground appearing in the finite certificate, in particular for $\Theta$;
\item the completed-block data $(C,R,\Theta\cup U)$ and its support function as in Lemma~\ref{lem:closed-subblock};
\item the algebraic operations on $M$, constants for all natural numbers, and the relevant conditional expectations.
\end{itemize}
Since the fixed certificate has finite complexity, only finitely many
previously constructed closure languages occur here; their union is still a
countable language.

For $S\subset\Omega$, let $\Hull_\Omega(S)$ be the hull of $S\cup\{\omega\}$ together with the structural parameters in this structure, and define
\[
   \cl_\Omega(S)=\Omega\cap\Hull_\Omega(S).
\]
Lemma~\ref{lem:skolem-hull} gives the size, monotonicity, and chain-continuity assertions, and the same hull calculation as in Lemma~\ref{lem:closed-subblock} gives idempotence.

Assume $S=\cl_\Omega(S)$ and put $\mathcal E=\Hull_\Omega(S)$.  Then $S=\Omega\cap\mathcal E$.  Let
\[
   S^- = S\cap\Omega^-,
   \qquad
   U_S=S\cap U,
   \qquad
   \Theta_S=\Theta\cap\mathcal E.
\]
Because the Skolem functions defining \(\cl_{\Omega^-}\) are included in the
present language, \(S^-\) is \(\cl_{\Omega^-}\)-closed.  Indeed,
\[
   \cl_{\Omega^-}(S^-)
   \subset \Omega^-\cap\mathcal E
   =
   S^-,
\]
and the reverse inclusion follows from the definition of a closure operation.
Hence \(S^-\) is certified by the induction hypothesis.  Let
\(\cl_\Theta\) be a strong closure operation for the certified background
\(\Theta\), supplied by the induction hypothesis for the finite certificate
of \(\Theta\).  The Skolem functions defining \(\cl_\Theta\) have been
included in the present language.  Therefore
\[
   \cl_\Theta(\Theta_S)
   \subset \Theta\cap\mathcal E
   =
   \Theta_S,
\]
and the reverse inclusion follows from extensivity of the closure operation.
Thus \(\Theta_S\) is \(\cl_\Theta\)-closed, and hence is a certified
subbackground.  Since \(\Theta\subset\Omega^-\), we also have
\(\Theta_S\subset S^-\).  Because the Skolem functions for the completed block
\(\mathfrak B=(C,R,\Theta\cup U)\) are included in the present language,
\[
   T_S=(\Theta\cup U)\cap\mathcal E=\Theta_S\cup U_S
\]
is closed for the block closure from Lemma~\ref{lem:closed-subblock}.  Indeed,
\[
   \cl_{\mathfrak B}(T_S)
   \subset
   (\Theta\cup U)\cap\mathcal E
   =
   T_S,
\]
and the reverse inclusion is automatic. Applying that lemma to this closed set gives von Neumann algebras $C_S\subset C$ and $R_S\subset R$ such that
\[
   L^2(R_S)_{\sa}=L^2(C_S)_{\sa}\oplus\overline{\spanop_{\RR}}(\Theta_S\cup U_S).
\]
Thus $U_S$ is a layer over the certified subbackground $\Theta_S\subset S^-$.  It follows that
\[
   S=S^-\dot\cup U_S
\]
is finite-layer certified.

For the final compatibility assertion, construct \(\cl_{\Omega^+}\) by using
a Skolem language which contains all Skolem functions used for \(\cl_\Omega\)
and the completed-block data defining the new layer \(U\) over \(\Theta\).
Let \(\Xi=\cl_{\Omega^+}(\Xi)\) and
\(\mathcal E=\Hull_{\Omega^+}(\Xi)\).  Then
\[
   \Xi=\Omega^+\cap\mathcal E,
   \qquad
   \Xi\cap\Omega=\Omega\cap\mathcal E.
\]
Since \(\mathcal E\) is closed under the finitary Skolem functions defining
\(\cl_\Omega\),
\[
   \cl_\Omega(\Xi\cap\Omega)\subset \Omega\cap\mathcal E=\Xi\cap\Omega.
\]
The reverse inclusion follows from extensivity, and therefore
\[
   \Xi\cap\Omega=\cl_\Omega(\Xi\cap\Omega).
\]

We also need the old support of the new layer to be certified.  Since
\(\Theta=\cl_\Omega(\Theta)\) and \(\mathcal E\) is closed under the Skolem
functions defining \(\cl_\Omega\), we have
\[
   \cl_\Omega(\Theta\cap\mathcal E)
   \subset
   \cl_\Omega(\Theta)\cap\mathcal E
   =
   \Theta\cap\mathcal E.
\]
Again the reverse inclusion follows from extensivity, so
\[
   \Theta\cap\mathcal E
   =
   \cl_\Omega(\Theta\cap\mathcal E).
\]
Thus \(\Theta\cap\mathcal E\) is a certified subbackground by the closure
assertions already proved for \(\Omega\).

Finally, apply Lemma~\ref{lem:closed-subblock} to the completed block
\((C,R,\Theta\cup U)\).  Since the corresponding block-closure Skolem
functions are included in the language,
\[
   (\Theta\cup U)\cap\mathcal E
   =
   (\Theta\cap\mathcal E)\cup(\Xi\cap U)
\]
is closed for that block closure.  Hence there are von Neumann algebras
\(C_\Xi\subset C\) and \(R_\Xi\subset R\) such that
\[
   L^2(R_\Xi)_{\sa}
   =
   L^2(C_\Xi)_{\sa}
   \oplus
   \overline{\spanop_{\RR}}\bigl((\Theta\cap\mathcal E)\cup(\Xi\cap U)\bigr).
\]
Thus \(\Xi\cap U\) is a certified layer over the certified subbackground
\(\Theta\cap\mathcal E\).  Since
\[
   \Theta\cap\mathcal E\subset \Xi\cap\Omega
\]
and \(\Xi\cap\Omega\) is \(\cl_\Omega\)-closed, the closure assertions already
proved for \(\Omega\) show that \(\Xi\cap\Omega\) is finite-layer certified.
Therefore
\[
   \Xi=(\Xi\cap\Omega)\dot\cup(\Xi\cap U)
\]
is finite-layer certified by adjoining the layer \(\Xi\cap U\) over the
certified subbackground \(\Theta\cap\mathcal E\subset\Xi\cap\Omega\).
\end{proof}




\begin{cor}[Restriction and adjoining a layer]\label{cor:cert-restrict-adjoin}
Let $(\Omega,\mathfrak c_\Omega)$ be a finite-layer certified background, and
let $\cl_\Omega$ be a closure operation supplied by
Proposition~\ref{prop:strong-closure}.

\begin{enumerate}[label=(\roman*)]
\item If $S\subset\Omega$ satisfies $S=\cl_\Omega(S)$, then $S$ admits a
      finite-layer certificate $\mathfrak c_S$.  With this certificate,
      $P_{\mathfrak c_S}(x)$ is the Hilbert-space projection of $x$ onto
      $\overline{\spanop_{\RR}}S$ for every $x=x^*\in M$.

\item Suppose $S=\cl_\Omega(S)$ and $U\subset\mathcal S_0(M)$ is an
      orthonormal family orthogonal to $\Omega$.  Suppose further that there
      are von Neumann algebras $C\subset R\subset M$ such that
      \[
         L^2(R)_{\sa}
         =
         L^2(C)_{\sa}\oplus
         \overline{\spanop_{\RR}}(S\cup U).
      \]
      Then $\Omega\dot\cup U$ is a finite-layer certified background.  A
      certificate is obtained by taking the old certificate
      $\mathfrak c_\Omega$, the certificate $\mathfrak c_S$ from part
      \emph{(i)}, and adjoining the single layer $U$ over $S$.

\item The strong closure operation on $\Omega\dot\cup U$ may be chosen
      compatibly with $\cl_\Omega$ in the sense of
      Proposition~\ref{prop:strong-closure}.
\end{enumerate}
\end{cor}

\begin{proof}
Part (i) is exactly Proposition~\ref{prop:strong-closure}(vi), together with
Lemma~\ref{lem:certified-projection}.  Part (ii) follows directly from the
recursive definition of finite-layer certificates in
Definition~\ref{def:certified}.  Part (iii) is the compatibility assertion in
Proposition~\ref{prop:strong-closure}.
\end{proof}











\section{The countable certified extension}\label{sec:countable-extension}

We first prove the extension theorem in the countable case.  The proof uses a standard fair bookkeeping device.

\begin{lem}[Fair bookkeeping with countable batches]\label{lem:bookkeeping}
Suppose that a recursion is indexed by \(n\geq1\), and that at each stage
at most countably many new tasks may appear.  Then the tasks can be labelled
and scheduled so that every task which eventually appears is processed
infinitely often at times \(n_k\) satisfying
\[
   \sum_k\frac1{n_k}=\infty .
\]
\end{lem}

\begin{proof}
For \(j\in\{0,1,2,\ldots\}\), put
\[
   A_j=\{n\geq1:\nu_2(n)=j\},
\]
where \(\nu_2(n)\) denotes the largest \(\nu\) such that \(2^\nu \mid n\).
Then the sets \(A_j\) partition the positive integers, and
\[
   \sum_{n\in A_j}\frac1n=\infty
   \qquad(j=0,1,2,\ldots).
\]

Next partition the label set \(\{0,1,2,\ldots\}\) into pairwise disjoint
infinite subsets
\[
   \{0,1,2,\ldots\}=\bigsqcup_{m=0}^\infty L_m .
\]
The labels in \(L_m\) will be reserved for tasks which first appear at stage
\(m\).

At stage \(m\), the new batch of tasks is countable, so we assign those tasks
injectively to unused labels in the infinite set \(L_m\).  At time \(n\geq1\),
let \(j=\nu_2(n)\).  If a task has already been assigned label \(j\), we
process that task; otherwise we do nothing at time \(n\).

If a task first appears at stage \(m\) and receives label \(j\in L_m\), then
it is processed at all sufficiently late times in \(A_j\).  Removing finitely
many early terms from \(A_j\) does not change the divergence of the harmonic
sum.  Hence the visiting times \(n_k\) of this task satisfy
\[
   \sum_k\frac1{n_k}=\infty .
\]
\end{proof}

\begin{lem}[Countable certified extension]\label{lem:countable-extension}
Let $\kappa=\dens L^2(M)>\aleph_0$.  Let $N\subset M$ be a von Neumann subalgebra with $\dens L^2(N)<\kappa$.  Let $\Omega\subset\mathcal S_0(M)$ be a finite-layer certified background with $|\Omega|<\kappa$, $E_N(\Omega)=0$, and a fixed strong closure operation $\cl_\Omega$ as in Proposition~\ref{prop:strong-closure}.  Let $S\subset\Omega$, $X\subset M_{\sa}$, and $Y\subset N_{\sa}$ be countable.  Then there are von Neumann algebras $C\subset N$ and $C\subset R\subset M$, a $\cl_\Omega$-closed set $\Omega_0\subset\Omega$, and an orthonormal family $U\subset\mathcal S_0(M)$ such that
\[
   S\subset\Omega_0,
   \qquad X\subset R,
   \qquad Y\subset C,
\]
\[
   |\Omega_0|,\ |U|,\ \dens L^2(R)\leq\aleph_0,
\]
\[
   E_N(U)=0,
   \qquad U\perp\Omega,
\]
and
\[
   L^2(R)_{\sa}=L^2(C)_{\sa}\oplus\overline{\spanop_{\RR}}(\Omega_0\cup U).
\]
Consequently $U$ is a layer over $\Omega_0$.
\end{lem}

\begin{proof}
We run a countable Henkin construction.  Maintain countable sets of parameters $\mathcal C\subset N_{\sa}$, $\Omega_*\subset\Omega$, and a finite or countable orthonormal family of already chosen symmetries $U_*$.  Initially
\[
   \mathcal C=Y,
   \qquad
   \Omega_* = \cl_\Omega(S),
   \qquad
   U_* =\varnothing.
\]
We implement the task list in countable batches, using
Lemma~\ref{lem:bookkeeping}.  


By a formal \(*\)-polynomial word we mean a noncommutative \(*\)-polynomial
with coefficients in \(\mathbb Q+i\mathbb Q\) in finitely many parameters
which have appeared by the relevant stage.  For such a word \(w\), we use the
self-adjoint parts
\[
   \Re w=\frac{w+w^*}{2},
   \qquad
   \Im w=\frac{w-w^*}{2i}.
\]
An \emph{eventual self-adjoint word} is an evaluated element of \(M_{\sa}\)
which is equal to \(\Re w\) or \(\Im w\) for some such formal word \(w\), after
all parameters occurring in \(w\) have appeared at some finite stage of the
construction.


At stage \(0\), declare all elements of the
initial parameter set
\[
   \mathcal C\cup X\cup\Omega_*\cup U_*
\]
to have appeared.  Enumerate all formal \( * \)-polynomial words with
coefficients in \(\mathbb Q+i\mathbb Q\) using these appeared parameters,
take their real and imaginary self-adjoint parts, and regard the resulting
self-adjoint words as the first countable batch of tasks.

At the end of each later stage, after \(\mathcal C\), \(\Omega_*\), or
\(U_*\) has possibly been enlarged, declare all newly added parameters to
have appeared.  Then enumerate all formal \( * \)-polynomial words with
coefficients in \(\mathbb Q+i\mathbb Q\) using the appeared parameters, take
their real and imaginary self-adjoint parts, and add precisely those
self-adjoint words which have not previously been labelled as the next
countable batch of tasks.

By Lemma~\ref{lem:bookkeeping}, the labels can be assigned so that every
eventual self-adjoint word is processed infinitely many times at visiting
times \(n_k\) satisfying
\[
   \sum_k\frac1{n_k}=\infty .
\]
At a stage whose scheduled label has not yet been assigned, we perform no
processing and continue to the next stage.  Otherwise the scheduled label
corresponds to a bounded self-adjoint word over the appeared parameter set,
and we process that word as follows.

Thus suppose that the current scheduled task is represented by a bounded self-adjoint word \(y\in M_{\sa}\). Define
\[
   b(y)=y-E_N(y)-P_\Omega(y).
\]
This is a bounded self-adjoint element of \(M\): indeed,
\(E_N(y)\in N_{\sa}\subset M_{\sa}\), and
\(P_\Omega(y)\in M_{\sa}\) by Lemma~\ref{lem:certified-projection}.

Moreover,
\[
   E_N(P_\Omega(y))=0,
\]
because \(P_\Omega(y)\) is an \(L^2\)-limit of finite real linear combinations
of elements of \(\Omega\), while \(E_N\) is \(L^2\)-continuous and vanishes on
\(\Omega\).  Hence
\[
   E_N(b(y))=0.
\]
Also
\[
   b(y)\perp\overline{\spanop_{\RR}}\Omega.
\]
Indeed, \(y-P_\Omega(y)\) is orthogonal to
\(\overline{\spanop_{\RR}}\Omega\) by definition of \(P_\Omega\), and
\(E_N(y)\perp\Omega\) because \(E_N(y)\in N_{\sa}\) and \(E_N(\Omega)=0\).

Add $E_N(y)$ to $\mathcal C$, add the countable set $\supp_\Omega(P_\Omega(y))$ to $\Omega_*$, and then replace $\Omega_*$ by $\cl_\Omega(\Omega_*)$.  This keeps $\Omega_*$ countable by Proposition~\ref{prop:strong-closure}.

At this finite stage, only finitely many symmetries have been chosen; write
them as \(u_1,\ldots,u_m\).  Since \(b(y)\) and the \(u_i\)'s are
self-adjoint, the real Hilbert-space projection is
\[
   P_{\spanop_{\RR}\{u_1,\ldots,u_m\}}b(y)
   =
   \sum_{i=1}^m \langle b(y),u_i\rangle_2 u_i,
\]
with real coefficients, and hence belongs to \(M_{\sa}\).  Put
\[
   r=b(y)-P_{\spanop_{\RR}\{u_1,\ldots,u_m\}}b(y).
\]
Thus \(r\in M_{\sa}\).  We record explicitly that \(r\) satisfies the
hypotheses needed for Lemma~\ref{lem:large-relative-norming}.  Since
\(E_N(b(y))=0\) and \(E_N(u_i)=0\) for \(1\leq i\leq m\), we have
\[
   E_N(r)=0.
\]
Moreover \(b(y)\perp\Omega\) and \(u_i\perp\Omega\) for all \(i\), hence
\[
   \tau(\omega r)=0\qquad(\omega\in\Omega).
\]
By construction \(r\) is the residual after orthogonal projection onto
\(\spanop_{\RR}\{u_1,\ldots,u_m\}\), so
\[
   \tau(u_i r)=0\qquad(1\leq i\leq m).
\]
Therefore \(r\) is orthogonal to
\[
   \Gamma=\Omega\cup\{u_1,\ldots,u_m\}.
\]
If $r=0$, choose no new symmetry at this stage.  If $r\neq0$, apply
Lemma~\ref{lem:large-relative-norming} with this constraint set \(\Gamma\).
This set has cardinal $<\kappa$.  We obtain $u_{m+1}\in\mathcal S_0(M)$ satisfying
\[
   E_N(u_{m+1})=0,
   \qquad
   u_{m+1}\perp\Omega,
   \qquad
   u_{m+1}\perp u_i\quad(1\leq i\leq m),
\]
and
\[
   \langle r,u_{m+1}\rangle_2
   =\frac{\|r\|_2^2}{\|r\|_\infty}.
\]
Add $u_{m+1}$ to $U_*$.

Let $U$ be the family of all chosen symmetries, let $\Omega_0$ be the union of the increasing sequence of values of $\Omega_*$, let $C=W^*(\mathcal C)$, and put
\[
   R=W^*(C,X,\Omega_0,U).
\]
The sets \(\Omega_0\) and \(U\) are countable.  Moreover, \(R\) is generated
by countably many bounded elements.  Hence, by
Lemma~\ref{lem:l2-cardinal-facts}\emph{(i)},
\[
   \dens L^2(R)\leq\aleph_0 .
\]
The set \(\Omega_0\) is \(\cl_\Omega\)-closed by the continuity of the closure
operation along the increasing sequence of values of \(\Omega_*\).

It remains to show that the announced decomposition holds.  Fix an eventual
self-adjoint word \(y\).  Once all parameters occurring in \(y\) have appeared,
the word \(y\) receives a fixed label.

For \(n\geq1\), let \(V_n(y)\) denote the real span of the symmetries which
have been chosen before stage \(n\).  Equivalently, if \(m_n\) symmetries have
been chosen before stage \(n\), then
\[
   V_n(y)=\spanop_{\RR}\{u_1,\ldots,u_{m_n}\}.
\]
Put
\[
   D_n(y)=\dist\bigl(b(y),V_n(y)\bigr).
\]
The sequence \(D_n(y)\) is decreasing, and its limit is
\[
   d=\dist\bigl(b(y),\overline{\spanop_{\RR}}U\bigr).
\]
We prove that \(d=0\).

Suppose, toward a contradiction, that \(d>0\).  Let
\[
   n_1<n_2<\cdots
\]
be the visiting times at which the fixed label of \(y\) is processed after
\(y\) has appeared.  By Lemma~\ref{lem:bookkeeping},
\[
   \sum_k\frac1{n_k}=\infty .
\]
At stage \(n_k\), the current residual is
\[
   r_k=b(y)-P_{V_{n_k}(y)}b(y).
\]
Hence
\[
   \|r_k\|_2=D_{n_k}(y)\geq d .
\]
In particular \(r_k\neq0\), so the construction chooses a new symmetry at
stage \(n_k\).

Since at most one new symmetry is chosen at each stage, we have \(m_{n_k}\leq
n_k\).  Also,
\[
\begin{aligned}
   \|r_k\|_\infty
   &\leq \|b(y)\|_\infty+
   \left\|P_{V_{n_k}(y)}b(y)\right\|_\infty       \\
   &\leq \|b(y)\|_\infty+
      \sum_{i=1}^{m_{n_k}}|\langle b(y),u_i\rangle_2|                         \\
   &\leq \|b(y)\|_\infty+\sqrt{m_{n_k}}\,\|b(y)\|_2                             \\
   &\leq K_y\sqrt{n_k},
\end{aligned}
\]
where
\[
   K_y=\|b(y)\|_\infty+\|b(y)\|_2 .
\]
If \(K_y=0\), then \(b(y)=0\) and there is nothing to prove.  Thus we may
assume \(K_y>0\).

At stage \(n_k\), let \(u\) be the new symmetry chosen by
Lemma~\ref{lem:large-relative-norming}.  Then \(u\perp V_{n_k}(y)\), and
\[
   \langle r_k,u\rangle_2=\frac{\|r_k\|_2^2}{\|r_k\|_\infty}.
\]
Therefore
\[
\begin{aligned}
   D_{n_k+1}(y)^2
   &\leq
   \dist\bigl(b(y),V_{n_k}(y)+\mathbb R u\bigr)^2  \\
   &=
   \|r_k\|_2^2-|\langle r_k,u\rangle_2|^2  \\
   &=
   D_{n_k}(y)^2-
   \frac{\|r_k\|_2^4}{\|r_k\|_\infty^2}.
\end{aligned}
\]
Since \(D_n(y)\) is decreasing, summing over \(k=1,\ldots,\ell\) gives
\[
   \sum_{k=1}^{\ell}
   \frac{\|r_k\|_2^4}{\|r_k\|_\infty^2}
   \leq D_{n_1}(y)^2 .
\]
Using \(\|r_k\|_2\ge d\) and
\(\|r_k\|_\infty\le K_y\sqrt{n_k}\), we get
\[
   \sum_{k=1}^{\ell}\frac{d^4}{K_y^2n_k}
   \leq D_{n_1}(y)^2 .
\]
Letting \(\ell\to\infty\) contradicts
\[
   \sum_k\frac1{n_k}=\infty .
\]
Therefore \(d=0\), and hence
\[
   b(y)\in\overline{\spanop_{\RR}}U
\]
for every eventual self-adjoint word \(y\).

For such $y$ we have
\[
   y=E_N(y)+P_\Omega(y)+b(y),
\]
where $E_N(y)\in L^2(C)_{\sa}$, $P_\Omega(y)\in\overline{\spanop_{\RR}}\Omega_0$ by the support bookkeeping, and $b(y)\in\overline{\spanop_{\RR}}U$. The eventual words are \(L^2\)-dense in \(L^2(R)_{\sa}\).  Let
\(\mathcal C_\infty\) be the final value of \(\mathcal C\), and put
\[
   \mathcal P_\infty=\mathcal C_\infty\cup X\cup\Omega_0\cup U .
\]
Then
\[
   C=W^*(\mathcal C_\infty),
   \qquad
   R=W^*(C,X,\Omega_0,U)=W^*(\mathcal P_\infty).
\]
Every element of \(\mathcal P_\infty\) appears at some finite stage of the
construction.  Hence every unital rational-coefficient \(*\)-word in
\(\mathcal P_\infty\), and its real and imaginary self-adjoint parts, is an
eventual word.  Let
\[
   \mathcal A_{\mathbb Q}
   =
   \operatorname{*-alg}_{\mathbb Q+i\mathbb Q}(\mathcal P_\infty).
\]
The self-adjoint part of \(\mathcal A_{\mathbb Q}\) is contained in the set
of eventual self-adjoint words.  Its norm closure is the \(C^*\)-algebra
generated by \(\mathcal P_\infty\), whose bicommutant is \(R\).  By
Kaplansky density, \((\mathcal A_{\mathbb Q})_{\sa}\) is \(L^2\)-dense in
\(L^2(R)_{\sa}\).  Therefore the eventual self-adjoint words are
\(L^2\)-dense in \(L^2(R)_{\sa}\).

The right-hand side below is closed, since
\(L^2(C)_{\sa}\) is orthogonal to
\(\overline{\spanop_{\RR}}(\Omega_0\cup U)\).  This orthogonality follows from
\(C\subset N\), \(E_N(\Omega_0)=0\), \(E_N(U)=0\), and \(U\perp\Omega\). Hence
\[
   L^2(R)_{\sa}\subset L^2(C)_{\sa}+\overline{\spanop_{\RR}}(\Omega_0\cup U).
\]
The reverse inclusion is clear from the definition of \(R\).  This proves the
completed-block decomposition.
\end{proof}

\section{The all-cardinal certified extension theorem}\label{sec:all-cardinal-extension}

We first record the elementary continuity fact used at limit stages.

\begin{lem}[Continuity of nested completed blocks]\label{lem:nested-limit}
Let $(C_i,R_i,T_i)_{i\in I}$ be an increasing chain of completed blocks, in the sense that $C_i\subset C_j$, $R_i\subset R_j$, and $T_i\subset T_j$ whenever $i\leq j$.  Put
\[
   C=W^*\Bigl(\bigcup_i C_i\Bigr),
   \qquad
   R=W^*\Bigl(\bigcup_i R_i\Bigr),
   \qquad
   T=\bigcup_iT_i.
\]
Then
\[
   L^2(R)_{\sa}=L^2(C)_{\sa}\oplus\overline{\spanop_{\RR}}T.
\]
\end{lem}

\begin{proof}
In a finite von Neumann algebra, if $R=W^*(\bigcup_iR_i)$ for an increasing family, then
\[
   L^2(R)=\overline{\bigcup_i L^2(R_i)}^{\|\cdot\|_2}.
\]
Indeed, the strong-$*$ closure of the algebraic union is $R$, and bounded strong-$*$ convergence implies $L^2$-convergence; Kaplansky density then gives the equality.  The same holds for $C$.

The subspace $L^2(C)_{\sa}$ is orthogonal to $\overline{\spanop_{\RR}}T$ because this orthogonality holds at every stage and passes to $L^2$-closures.  The closed sum contains each $L^2(R_i)_{\sa}$, hence contains the $L^2$-closure of their union, which is $L^2(R)_{\sa}$.  The reverse inclusion is immediate from $C\subset R$ and $T\subset R$.
\end{proof}

\begin{rem}\label{rem:no-regularity}
No regularity assumption on $\kappa$ is used.  We identify cardinals with
initial ordinals.  If $\lambda>\aleph_0$ is an initial cardinal and
$\alpha<\lambda$, then $|\alpha|<\lambda$, and hence
\[
   \max\{\aleph_0,|\alpha|\}<\lambda .
\]
The case $\lambda=\aleph_0$ is handled separately as the base case in
Proposition~\ref{prop:LBE}; in that case the displayed strict inequality with
$\lambda$ on the right is not used.

In all applications below, whenever $\lambda<\kappa$ and $\alpha<\lambda$, one
still has
\[
   \max\{\aleph_0,|\alpha|\}<\kappa .
\]
Indeed, if $\lambda=\aleph_0$, then the left-hand side is $\aleph_0<\kappa$;
if $\lambda>\aleph_0$, then the preceding paragraph gives
\(\max\{\aleph_0,|\alpha|\}<\lambda<\kappa\).  Thus the argument does not use
regularity of $\kappa$, and this remains true when $\kappa$ is singular.
\end{rem}

\begin{prop}[Certified extension theorem]\label{prop:LBE}
Let $\kappa=\dens L^2(M)>\aleph_0$ and let $\aleph_0\leq\lambda<\kappa$.  Let $N\subset M$ satisfy $\dens L^2(N)<\kappa$.  Let $\Omega\subset\mathcal S_0(M)$ be a finite-layer certified background with $|\Omega|<\kappa$, $E_N(\Omega)=0$, and a fixed strong closure operation $\cl_\Omega$ as in Proposition~\ref{prop:strong-closure}.  Let
\[
   S\subset\Omega,
   \qquad
   X\subset M_{\sa},
   \qquad
   Y\subset N_{\sa}
\]
with $|S|,|X|,|Y|\leq\lambda$.  Then there are von Neumann algebras $C\subset N$ and $C\subset R\subset M$, a $\cl_\Omega$-closed set $\Omega_0\subset\Omega$, and an orthonormal family $U\subset\mathcal S_0(M)$ such that
\[
   S\subset\Omega_0,
   \qquad X\subset R,
   \qquad Y\subset C,
\]
\[
   |\Omega_0|,\ |U|,\ \dens L^2(R)\leq\lambda,
\]
\[
   E_N(U)=0,
   \qquad U\perp\Omega,
\]
and
\[
   L^2(R)_{\sa}=L^2(C)_{\sa}\oplus\overline{\spanop_{\RR}}(\Omega_0\cup U).
\]
Thus $U$ is a layer over $\Omega_0$.
\end{prop}

\begin{proof}
We prove the proposition by induction on $\lambda$.  The case $\lambda=\aleph_0$ is Lemma~\ref{lem:countable-extension}.

Assume $\lambda>\aleph_0$ and that the proposition has been proved for all cardinals $\mu$ with $\aleph_0\leq\mu<\lambda$.  Enumerate the disjoint union of the data $S$, $X$, and $Y$ as $(z_\alpha)_{\alpha<\lambda}$, with repetitions or dummy entries if necessary, remembering the type of each genuine datum.  We first apply the countable case to the original background $\Omega$, with support $\cl_\Omega(\varnothing)$ and with $X=Y=\varnothing$.  This gives an initial completed state
\[
   (C_0,R_0,\Omega_0\cup U_0)
\]
with countable density, where $\Omega_0$ is $\cl_\Omega$-closed and $U_0$ is a layer over $\Omega_0$.  The recursion below starts from this state; at successor stages we enlarge the current state to absorb the next datum.

We recursively construct, for $\alpha<\lambda$, completed states
\[
   (C_\alpha,R_\alpha,\Omega_\alpha\cup U_\alpha)
\]
such that
\[
   C_\alpha\subset N,
   \quad C_\alpha\subset R_\alpha\subset M,
   \quad \Omega_\alpha\subset\Omega,
   \quad \Omega_\alpha=\cl_\Omega(\Omega_\alpha),
\]
$U_\alpha$ is a layer over $\Omega_\alpha$,
\[
   E_N(U_\alpha)=0,
   \quad U_\alpha\perp\Omega,
\]
and
\[
   \dens L^2(R_\alpha),\ |\Omega_\alpha|,\ |U_\alpha|
   \leq\max\{\aleph_0,|\alpha|\}.
\]
The states are required to be increasing in all coordinates.

Suppose the state at $\alpha$ is constructed.  Put
\[
   \lambda_\alpha=\max\{\aleph_0,|\alpha|\}<\lambda.
\]
Here \(\lambda\) is identified with its initial ordinal.  No regularity of
\(\lambda\), and hence no regularity of \(\kappa\), is used in this inequality;
see Remark~\ref{rem:no-regularity}.
By the induction invariant,
\[
   \Omega_\alpha=\cl_\Omega(\Omega_\alpha)
\]
and
\[
   L^2(R_\alpha)_{\sa}
   =
   L^2(C_\alpha)_{\sa}\oplus
   \overline{\spanop_{\RR}}(\Omega_\alpha\cup U_\alpha).
\]
Moreover $U_\alpha$ is orthonormal and orthogonal to $\Omega$. Therefore
Corollary~\ref{cor:cert-restrict-adjoin}, applied with
$S=\Omega_\alpha$ and $U=U_\alpha$, certifies the background
\[
   \Omega\dot\cup U_\alpha .
\]
This is a re-certification of the whole currently constructed family
\(U_\alpha\) as a single layer over \(\Omega_\alpha\); the finite certificate
depth is not accumulated along the transfinite recursion.  The same corollary
supplies a strong closure operation on
$\Omega\dot\cup U_\alpha$ compatible with the original operation
$\cl_\Omega$. Also
\[
   E_N(\Omega\dot\cup U_\alpha)=0
\]
and
\[
   |\Omega\dot\cup U_\alpha|<\kappa,
\]
because $|\Omega|<\kappa$ and
$|U_\alpha|\leq\lambda_\alpha<\lambda<\kappa$.

Choose self-adjoint \(L^2\)-dense sets
\[
   Y_\alpha\subset (C_\alpha)_{\sa},
   \qquad
   X_\alpha\subset (R_\alpha)_{\sa},
\]
of cardinality at most \(\lambda_\alpha\).  This is possible because
\[
   \dens L^2(C_\alpha)\leq \dens L^2(R_\alpha)\leq\lambda_\alpha
\]
and the bounded self-adjoint parts are \(L^2\)-dense in the corresponding
self-adjoint \(L^2\)-spaces.  By Lemma~\ref{lem:l2-cardinal-facts}(iv),
\[
   W^*(Y_\alpha)=C_\alpha,
   \qquad
   W^*(X_\alpha)=R_\alpha.
\]
Set
\[
   S'_\alpha=\Omega_\alpha\cup U_\alpha,
   \qquad X'_\alpha=X_\alpha,
   \qquad Y'_\alpha=Y_\alpha.
\]
If the datum $z_\alpha$ comes from the original set $S$, add it to $S'_\alpha$; if it comes from $X$, add it to $X'_\alpha$; and if it comes from $Y$, add it to $Y'_\alpha$.  Apply the induction hypothesis at cardinal $\lambda_\alpha$ to the background $\Omega\dot\cup U_\alpha$ with input $S'_\alpha,X'_\alpha,Y'_\alpha$.  We obtain a completed extension
\[
   (C_{\alpha+1},R_{\alpha+1},\Xi_{\alpha+1}\cup V_{\alpha+1}),
\]
where
\[
   \Xi_{\alpha+1}\subset\Omega\dot\cup U_\alpha
\]
is closed for the chosen closure operation on $\Omega\dot\cup U_\alpha$ and contains $\Omega_\alpha\cup U_\alpha$.  Since all of $U_\alpha$ is required support,
\[
   \Xi_{\alpha+1}=(\Xi_{\alpha+1}\cap\Omega)\dot\cup U_\alpha.
\]
Set
\[
   \Omega_{\alpha+1}=\Xi_{\alpha+1}\cap\Omega,
   \qquad
   U_{\alpha+1}=U_\alpha\cup V_{\alpha+1}.
\]
By the compatibility part of Proposition~\ref{prop:strong-closure},
\(\Omega_{\alpha+1}\) is \(\cl_\Omega\)-closed.  Moreover the new state
contains the old one.  Indeed, \(Y_\alpha\subset C_{\alpha+1}\) and
\(X_\alpha\subset R_{\alpha+1}\) by the choice of the input data for the
induction hypothesis; since
\[
   W^*(Y_\alpha)=C_\alpha,
   \qquad
   W^*(X_\alpha)=R_\alpha,
\]
we get
\[
   C_\alpha\subset C_{\alpha+1},
   \qquad
   R_\alpha\subset R_{\alpha+1}.
\]
Also \(\Omega_\alpha\subset\Omega_{\alpha+1}\) and
\(U_\alpha\subset U_{\alpha+1}\) by construction.  Thus the completed states
are increasing in all coordinates.

The completed-block decomposition for the output says exactly that
\[
   L^2(R_{\alpha+1})_{\sa}=L^2(C_{\alpha+1})_{\sa}\oplus
   \overline{\spanop_{\RR}}(\Omega_{\alpha+1}\cup U_{\alpha+1}).
\]
The cardinal bounds are $\leq\lambda_\alpha\leq\max\{\aleph_0,|\alpha+1|\}$.  Thus the successor step is complete.

At a limit ordinal $\delta<\lambda$, define
\[
   C_\delta=W^*\Bigl(\bigcup_{\alpha<\delta}C_\alpha\Bigr),
   \qquad
   R_\delta=W^*\Bigl(\bigcup_{\alpha<\delta}R_\alpha\Bigr),
\]
\[
   \Omega_\delta=\bigcup_{\alpha<\delta}\Omega_\alpha,
   \qquad
   U_\delta=\bigcup_{\alpha<\delta}U_\alpha.
\]
Since \(\delta<\lambda\) and \(\lambda\) is an initial cardinal,
\(|\delta|<\lambda\).  The cardinal estimates are as follows.  For each
\(\alpha<\delta\), choose an \(L^2\)-dense subset
\(D_\alpha\subset L^2(R_\alpha)\) of cardinal at most
\(\max\{\aleph_0,|\alpha|\}\).  Then
\[
   \left|\bigcup_{\alpha<\delta}D_\alpha\right|
   \leq
   |\delta|\cdot\max\{\aleph_0,|\delta|\}
   =
   \max\{\aleph_0,|\delta|\}.
\]
By Lemma~\ref{lem:l2-cardinal-facts}(ii), this union is \(L^2\)-dense in
\(L^2(R_\delta)\).  The same cardinal calculation gives
\[
   |\Omega_\delta|,\ |U_\delta|
   \leq\max\{\aleph_0,|\delta|\}.
\]
Hence
\[
   \dens L^2(R_\delta),\ |\Omega_\delta|,\ |U_\delta|
   \leq\max\{\aleph_0,|\delta|\}<\lambda.
\]
The continuity of $\cl_\Omega$ from Proposition~\ref{prop:strong-closure} shows that $\Omega_\delta$ is $\cl_\Omega$-closed.  Lemma~\ref{lem:nested-limit} applied to the nested completed blocks gives
\[
   L^2(R_\delta)_{\sa}=L^2(C_\delta)_{\sa}\oplus
   \overline{\spanop_{\RR}}(\Omega_\delta\cup U_\delta).
\]
Thus $U_\delta$ is a layer over $\Omega_\delta$.

Finally put
\[
   C=W^*\Bigl(\bigcup_{\alpha<\lambda}C_\alpha\Bigr),
   \qquad
   R=W^*\Bigl(\bigcup_{\alpha<\lambda}R_\alpha\Bigr),
\]
\[
   \Omega_{\mathrm{out}}=\bigcup_{\alpha<\lambda}\Omega_\alpha,
   \qquad
   U=\bigcup_{\alpha<\lambda}U_\alpha.
\]
The continuity of $\cl_\Omega$ gives that $\Omega_{\mathrm{out}}$ is $\cl_\Omega$-closed, and Lemma~\ref{lem:nested-limit} gives
\[
   L^2(R)_{\sa}=L^2(C)_{\sa}\oplus
   \overline{\spanop_{\RR}}(\Omega_{\mathrm{out}}\cup U).
\]
All data from $S$, $X$, and $Y$ were absorbed during the recursion.  The
cardinal bounds for \(\Omega_{\mathrm{out}}\) and \(U\) follow from taking a
union of \(\lambda\) many sets each of cardinal at most \(\lambda\):
\[
   |\Omega_{\mathrm{out}}|,\ |U|\leq \lambda .
\]

It remains only to record the density bound for \(R\).  For each
\(\alpha<\lambda\), choose an \(L^2\)-dense set
\(D_\alpha\subset L^2(R_\alpha)\) with \(|D_\alpha|\leq\lambda\).  Then
\[
   \left|\bigcup_{\alpha<\lambda}D_\alpha\right|
   \leq \lambda\cdot\lambda
   =\lambda .
\]
By Lemma~\ref{lem:l2-cardinal-facts}\emph{(ii)}, the union
\(\bigcup_{\alpha<\lambda}D_\alpha\) is \(L^2\)-dense in \(L^2(R)\).  Hence
\[
   \dens L^2(R)\leq\lambda .
\]
Thus the conclusion holds with \(\Omega_0=\Omega_{\mathrm{out}}\).
\end{proof}

\section{Relative extension and the nonseparable theorem}\label{sec:nonseparable-main}

\begin{prop}[Relative extension property]\label{prop:RE}
Let $\kappa=\dens L^2(M)>\aleph_0$ and let $\aleph_0\leq\lambda<\kappa$.  Let $N\subset M$ be a von Neumann subalgebra with $\dens L^2(N)\leq\lambda$.  Suppose $H_0(N)$ has an orthonormal basis $B_N\subset\mathcal S_0(N)$.  If $X\subset M_{\sa}$ and $|X|\leq\lambda$, then there are a von Neumann algebra $P$ with $N\cup X\subset P\subset M$ and an orthonormal family \(U\subset\mathcal S_0(P)\) such that
\[
   \dens L^2(P)\leq\lambda,
\]
\[
   U\perp L^2(N)_{\sa},
\]
and
\[
   B_N\cup U
\]
is an orthonormal basis of $H_0(P)$.
\end{prop}

\begin{proof}
Apply Proposition~\ref{prop:LBE} with the empty background
$\Omega=\varnothing$, $S=\varnothing$, and with
$Y\subset N_{\sa}$ a self-adjoint $L^2$-dense set of cardinal
$\leq\lambda$.  We obtain $C\subset N$, $C\subset R\subset M$, and
$U\subset\mathcal S_0(M)$ such that $Y\subset C$, $X\subset R$, and
\[
   L^2(R)_{\sa}=L^2(C)_{\sa}\oplus\overline{\spanop_{\RR}}U.
\]
Since $Y$ is $L^2$-dense in $L^2(N)_{\sa}$,
Lemma~\ref{lem:l2-cardinal-facts}\emph{(iv)} gives $W^*(Y)=N$.
Hence $C=N$.  Put \(P=R\).  Then \(N=C\subset R=P\) and \(X\subset P\), and
\[
   L^2(P)_{\sa}=L^2(N)_{\sa}\oplus\overline{\spanop_{\RR}}U.
\]
In particular,
\[
   U\perp L^2(N)_{\sa}.
\]
Moreover \(U\subset P\).  Indeed, each \(u\in U\) belongs to \(M\), and its
\(L^2\)-vector belongs to \(L^2(P)\) by the displayed decomposition.  Hence
Lemma~\ref{lem:l2-cardinal-facts}\emph{(iii)} gives \(u\in P\).  Since each
\(u\in U\) is a trace-zero symmetry in \(M\), it is also a trace-zero symmetry
in \(P\).  Thus \(U\subset\mathcal S_0(P)\).

Since \(U\subset\mathcal S_0(P)\), every element of \(U\) is orthogonal to
\(\one\). Taking the orthogonal complement of \(\RR\one\) in the displayed
decomposition gives
\[
   H_0(P)=H_0(N)\oplus\overline{\spanop_{\RR}}U .
\]
Since \(B_N\) is an orthonormal basis of \(H_0(N)\), the family
\(B_N\cup U\) is an orthonormal basis of \(H_0(P)\).
\end{proof}

\begin{lem}[A separable abelian starting algebra]\label{lem:start}
Every $\mathrm{II}_1$ factor contains a separable diffuse abelian von Neumann subalgebra $A_0$.  Moreover $H_0(A_0)$ has an orthonormal basis contained in $\mathcal S_0(A_0)$.
\end{lem}

\begin{proof}
We construct a dyadic system of projections.  Start with $p_\varnothing=1$.
By the comparison and dimension theory of projections in a finite factor
\cite[Ch.~6]{KR97II}, every nonzero projection in $M$ can be split into two
orthogonal subprojections of equal trace.  Recursively, for every finite
binary string $\sigma\in\{0,1\}^{<\mathbb N}$, choose orthogonal projections
$p_{\sigma 0},p_{\sigma 1}\leq p_\sigma$ such that
\[
   p_{\sigma 0}+p_{\sigma 1}=p_\sigma,
   \qquad
   \tau(p_{\sigma 0})=\tau(p_{\sigma 1})=\frac12\tau(p_\sigma).
\]
The projections in this dyadic tree commute pairwise, since any two of them
are either orthogonal or one is dominated by the other. Let \(A_0\) be the abelian von Neumann algebra generated by all projections
\(p_\sigma\). For \(n\geq0\), put
\[
   F_n=\spanop\{p_\sigma:|\sigma|=n\}.
\]
Then \(F_n\subset F_{n+1}\), and
\[
   A_0=W^*\Bigl(\bigcup_{n\geq0}F_n\Bigr).
\]
Hence, by Lemma~\ref{lem:l2-cardinal-facts}\emph{(ii)},
\[
   L^2(A_0)=
   \overline{\bigcup_{n\geq0}L^2(F_n)}^{\|\cdot\|_2},
\]
so \(L^2(A_0)\) is separable.

We claim that \(A_0\) is diffuse.  Let \(0\neq q\in A_0\) be a projection.
Choose \(n\) with \(2^{-n}<\tau(q)\).  Since
\(\sum_{|\sigma|=n}p_\sigma=\one\), there is some \(\sigma\in\{0,1\}^n\)
such that \(qp_\sigma\neq0\).  If \(qp_\sigma=q\), then \(q\le p_\sigma\),
which would imply
\[
   \tau(q)\leq\tau(p_\sigma)=2^{-n},
\]
a contradiction.  Thus
\[
   0<qp_\sigma<q.
\]
Since \(A_0\) is abelian, \(qp_\sigma\) is a projection in \(A_0\).  Hence
\(q\) is not minimal, and \(A_0\) is diffuse.

For every finite set \(F\subset\mathbb N\), choose \(n\) with
\(F\subset\{1,\ldots,n\}\), and define
\[
   w_F
   =
   \sum_{\sigma\in\{0,1\}^n}
   (-1)^{\sum_{j\in F}\sigma_j}p_\sigma .
\]
This definition is independent of the choice of \(n\).  Each \(w_F\) is a
self-adjoint unitary in \(A_0\), \(w_\varnothing=\one\), and
\(\tau(w_F)=0\) whenever \(F\neq\varnothing\).  For each \(n\), the finite
family
\[
   \{w_F:F\subset\{1,\ldots,n\}\}
\]
is the usual Walsh orthonormal basis of \(L^2(F_n)_{\sa}\). Therefore
\[
   \{w_F:F\subset\mathbb N,\ |F|<\infty\}
\]
is an orthonormal basis of \(L^2(A_0)_{\sa}\), and the nonconstant Walsh
symmetries
\[
   \{w_F:0<|F|<\infty\}
\]
form an orthonormal basis of \(H_0(A_0)\) contained in
\(\mathcal S_0(A_0)\).
\end{proof}

\begin{thm}[Nonseparable case]\label{thm:nonseparable-main}
Let $M$ be a $\mathrm{II}_1$ factor with faithful normal tracial state $\tau$,
and set
\(\kappa=\dens L^2(M,\tau)\).  If
\[
   \kappa>\aleph_0,
\]
then there is a family \((s_i)_{i\in I}\) of self-adjoint unitaries in $M$, with
\(|I|=\kappa\), such that \(\{\widehat{s_i}:i\in I\}\) is a Hilbert
orthonormal basis of $L^2(M,\tau)$.  Consequently $L^2(M,\tau)$ has a Hilbert
orthonormal basis consisting of trace vectors for the left action of $M$.
\end{thm}

\begin{proof}
Let
\[
   \kappa=\dens L^2(M,\tau)>\aleph_0.
\]
The decompositions
\[
   L^2(M)_{\sa}=H_0(M)\oplus\RR\one,
   \qquad
   L^2(M)=L^2(M)_{\sa}+iL^2(M)_{\sa}
\]
show that $\dens H_0(M)=\kappa$. The bounded trace-zero self-adjoint elements $M_{\sa}\cap H_0(M)$ are
$L^2$-dense in $H_0(M)$.  Indeed, if $\xi\in H_0(M)$, then by definition of
$L^2(M)_{\sa}$ there are $x_i=x_i^*\in M$ with
$\|x_i-\xi\|_2\to0$.  Since
\[
   \tau(x_i)=\langle x_i,1\rangle_2\longrightarrow
   \langle \xi,1\rangle_2=0,
\]
we have
\[
   x_i-\tau(x_i)1\in M_{\sa}\cap H_0(M)
\]
and
\[
   \|x_i-\tau(x_i)1-\xi\|_2\to0.
\] 
Choose an $L^2$-dense family
\[
   \{a_\alpha:\alpha<\kappa\}\subset M_{\sa}\cap H_0(M).
\]

By Lemma~\ref{lem:start}, choose a separable diffuse abelian von Neumann subalgebra $N_0\subset M$ and an orthonormal basis $B_0\subset\mathcal S_0(N_0)$ of $H_0(N_0)$.

We recursively construct von Neumann algebras $N_\alpha\subset M$ and orthonormal families $B_\alpha\subset\mathcal S_0(M)$, for $\alpha<\kappa$, such that:
\begin{enumerate}[label=(\roman*)]
\item \(B_\alpha\subset\mathcal S_0(N_\alpha)\) is an orthonormal basis of \(H_0(N_\alpha)\);
\item $a_\alpha\in N_{\alpha+1}$;
\item if $\alpha<\beta$, then $N_\alpha\subset N_\beta$ and $B_\alpha\subset B_\beta$;
\item $\dens L^2(N_\alpha)\leq\max\{\aleph_0,|\alpha|\}$.
\end{enumerate}
The initial pair is $(N_0,B_0)$.

Suppose $N_\alpha,B_\alpha$ have been constructed.  Put
\[
   \lambda_\alpha=\max\{\aleph_0,|\alpha|\}.
\]
Since $\alpha<\kappa$ and $\kappa$ is identified with its initial ordinal, $\lambda_\alpha<\kappa$.  Proposition~\ref{prop:RE}, applied to \(N=N_\alpha\) and
\(X=\{a_\alpha\}\), gives a von Neumann algebra \(N_{\alpha+1}\) containing
\(N_\alpha\cup\{a_\alpha\}\) and an orthonormal family
\(U_\alpha\subset\mathcal S_0(N_{\alpha+1})\) such that
\[
   B_\alpha\cup U_\alpha
\]
is an orthonormal basis of \(H_0(N_{\alpha+1})\).  Set
\[
   B_{\alpha+1}=B_\alpha\cup U_\alpha .
\]
The density bound is also supplied by Proposition~\ref{prop:RE}.

At a limit ordinal $\delta<\kappa$, set
\[
   N_\delta=W^*\Bigl(\bigcup_{\alpha<\delta}N_\alpha\Bigr),
   \qquad
   B_\delta=\bigcup_{\alpha<\delta}B_\alpha.
\]
In a finite von Neumann algebra,
\[
   L^2(N_\delta)=
   \overline{\bigcup_{\alpha<\delta}L^2(N_\alpha)}^{\|\cdot\|_2},
\]
by the same Kaplansky-density argument as in Lemma~\ref{lem:nested-limit}.
We claim that
\[
   H_0(N_\delta)
   =
   \overline{\bigcup_{\alpha<\delta}H_0(N_\alpha)}^{\|\cdot\|_2}.
\]
Let \(\xi\in H_0(N_\delta)\) and let \(\varepsilon>0\).  By the displayed
\(L^2\)-continuity and by taking self-adjoint parts, there are
\(\alpha<\delta\) and \(x=x^*\in N_\alpha\) such that
\[
   \|x-\xi\|_2<\varepsilon.
\]
Since \(\xi\perp\one\) and \(\|\one\|_2=1\),
\[
   |\tau(x)|
   =
   |\langle x-\xi,\one\rangle_2|
   \leq \|x-\xi\|_2
   <\varepsilon.
\]
Therefore \(x-\tau(x)\one\in H_0(N_\alpha)\), and
\[
   \|x-\tau(x)\one-\xi\|_2
   \leq
   \|x-\xi\|_2+|\tau(x)|\,\|\one\|_2
   <2\varepsilon.
\]
This proves the claim.
Since the bases are nested, $B_\delta$ is an orthonormal basis of $H_0(N_\delta)$.  Also
\[
   \dens L^2(N_\delta)\leq\max\{\aleph_0,|\delta|\}<\kappa.
\]

Put
\[
   B=\bigcup_{\alpha<\kappa}B_\alpha .
\]
Since the \(B_\alpha\)'s are nested orthonormal families, \(B\) is orthonormal, and \(B\subset\mathcal S_0(M)\).
For every $\alpha<\kappa$, the element $a_\alpha$ belongs to $N_{\alpha+1}$ and has trace zero; hence
\[
   a_\alpha\in H_0(N_{\alpha+1})\subset\overline{\spanop_{\RR}}B.
\]
Since the family $(a_\alpha)_{\alpha<\kappa}$ is $L^2$-dense in $H_0(M)$, the orthonormal family $B$ is a Hilbert orthonormal basis of $H_0(M)$.  Hence
\(|B|=\dens H_0(M)=\kappa\), because the density character of an infinite
Hilbert space equals the cardinality of any Hilbert orthonormal basis.  By
Lemma~\ref{lem:trace-zero-reduction}, $\{\widehat{\one}\}\cup\{\widehat{s}:s\in B\}$ is a Hilbert orthonormal
basis of $L^2(M,\tau)$ consisting of self-adjoint unitaries.  Since
\(\kappa>\aleph_0\), this basis has cardinality \(\kappa\).  Each of these
vectors is a trace vector, again by Lemma~\ref{lem:trace-zero-reduction}.
\end{proof}

\begin{rem}\label{rem:nonsep-density-gap}
The nonseparable argument uses the strict density gap
$\dens L^2(N)<\dens L^2(M)$ in Lemma~\ref{lem:large-relative-norming}.
This is why the proof is organized as a transfinite chain of subalgebras of
smaller density character.  The separable case is handled in \cite{HTZ26}
by a finite-dimensional norming argument and is not used in the present proof.
\end{rem}

\section{Consequence together with the separable case}
\label{sec:consequence}

\begin{cor}[Kadison's trace-vector basis problem]\label{cor:kadison-full}
Let $M$ be a type $\mathrm{II}_1$ factor with faithful normal tracial state
$\tau$.
Then $L^2(M,\tau)$ admits a Hilbert orthonormal basis whose vectors are
represented by self-adjoint unitaries in $M$.  Consequently $L^2(M,\tau)$ admits a Hilbert
orthonormal basis consisting of trace vectors for the left action of $M$.
\end{cor}

\begin{proof}
If $\dens L^2(M,\tau)=\aleph_0$, this is the separable case proved in
\cite{HTZ26}.  If $\dens L^2(M,\tau)>\aleph_0$, this is
Theorem~\ref{thm:nonseparable-main}.  In both cases the basis vectors are
represented by self-adjoint unitaries in $M$.  For every self-adjoint unitary
$s\in M$ and every $x\in M$,
\[
   \langle x\widehat s,\widehat s\rangle_2
   =\tau(sxs)
   =\tau(xss)
   =\tau(x),
\]
so each basis vector is a trace vector.
\end{proof}


\section*{Acknowledgments} 
Quanyu Tang thank Changying Ding and Nik Weaver for responding to inquiries concerning this candidate manuscript. Neither of them was able to review or verify the mathematical argument. Their mention here should not be understood as an endorsement of the manuscript or of any of its claims.

\begin{thebibliography}{KR97II}

\bibitem[AW03]{AW03}
C.~A. Akemann and N.~Weaver,
\emph{Automatic convexity},
J. Convex Anal. \textbf{10} (2003), no.~1, 275--284.
\url{https://www.heldermann-verlag.de/jca/jca10/jca0339.pdf}

\bibitem[CK90]{CK90}
C.~C. Chang and H.~J. Keisler,
\emph{Model Theory}, 3rd ed.,
Studies in Logic and the Foundations of Mathematics, vol.~73, North-Holland,
Amsterdam, 1990.

\bibitem[Con90]{Con90}
J.~B. Conway,
\emph{A Course in Functional Analysis},
2nd ed., Graduate Texts in Mathematics, vol.~96, Springer, New York, 1990.



\bibitem[Ge03]{Ge03}
L.~M. Ge,
\emph{On ``Problems on von Neumann Algebras by R. Kadison, 1967''},
Acta Math. Sin. (Engl. Ser.) \textbf{19} (2003), no.~3, 619--624.
\doi{10.1007/s10114-003-0279-x}

\bibitem[HTZ26]{HTZ26}
Y.~He, Q.~Tang, and T.~Zhang,
\emph{The separable case of Kadison's problem on orthonormal bases of
unitaries for type $\mathrm{II}_1$ factors},
arXiv:2605.15006v2 [math.OA], 2026.
\doi{10.48550/arXiv.2605.15006}

\bibitem[Jech03]{Jech03}
T.~Jech,
\emph{Set Theory: The Third Millennium Edition, Revised and Expanded},
Springer Monographs in Mathematics, Springer, Berlin, 2003.
\doi{10.1007/3-540-44761-X}

\bibitem[Kad67]{Kad67}
R.~V. Kadison,
\emph{Problems on von Neumann algebras},
The Baton Rouge Conference on Operator Algebras, Baton Rouge, LA, 1967,
unpublished.

\bibitem[KR97a]{KR97I}
R.~V. Kadison and J.~R. Ringrose,
\emph{Fundamentals of the Theory of Operator Algebras. Vol.~I: Elementary Theory},
Graduate Studies in Mathematics, vol.~15, American Mathematical Society,
Providence, RI, 1997.

\bibitem[KR97b]{KR97II}
R.~V. Kadison and J.~R. Ringrose,
\emph{Fundamentals of the Theory of Operator Algebras. Vol.~II: Advanced Theory},
Graduate Studies in Mathematics, vol.~16, American Mathematical Society,
Providence, RI, 1997.

\bibitem[Pet20]{Pet20}
J.~Peterson,
\emph{Open problems in operator algebras},
webpage, created October 9, 2020; last updated May 15, 2026.
Available at \url{https://www.math.uwaterloo.ca/~j37peter/problems.html}.

\bibitem[Tak72]{Tak72}
M.~Takesaki,
\emph{Conditional expectations in von Neumann algebras},
J. Functional Analysis \textbf{9} (1972), 306--321.
\doi{10.1016/0022-1236(72)90004-3}

\bibitem[Tak02]{Tak02}
M.~Takesaki,
\emph{Theory of Operator Algebras I},
Encyclopaedia of Mathematical Sciences, vol.~124, Springer, Berlin, 2002.

\end{thebibliography}

\end{document}
